% Generated by code/assemble_article.py; edit the sections/ fragments.

% ----- sections/preamble.tex -----
\documentclass[11pt,a4paper]{article}
\usepackage[T1]{fontenc}
\usepackage[utf8]{inputenc}
\usepackage{lmodern}
\usepackage[margin=25mm,headheight=15pt]{geometry}
\usepackage{amsmath,amssymb,amsthm,mathtools}
\usepackage{microtype}
\usepackage{graphicx}
\usepackage{booktabs,array,longtable}
\usepackage{enumitem}
\usepackage{xcolor}
\usepackage{xurl}
\usepackage{titlesec}
\usepackage{fancyhdr}
\definecolor{Ink}{HTML}{1B263B}
\definecolor{Teal}{HTML}{176C70}
\definecolor{Muted}{HTML}{566273}
\usepackage[colorlinks=true,linkcolor=Ink,citecolor=Teal,urlcolor=Teal,
  bookmarksnumbered=true,pdfencoding=auto]{hyperref}
\hypersetup{
 pdftitle={Real Zeros and Exact Relation Spaces in the Polylogarithm--Stieltjes Programme},
 pdfauthor={Research draft prepared with OpenAI ChatGPT for Vladimir Reshetnikov},
 pdfsubject={Stieltjes derivatives, weighted distribution modules, polylogarithm reductions, and corrections to ProveIt drafts},
 pdfkeywords={generalized Stieltjes constants, Hurwitz zeta, real zeros, Appell polynomials, distribution relations, polylogarithms}
}
\DeclareMathOperator{\Li}{Li}
\newtheorem{theorem}{Theorem}[section]
\newtheorem{lemma}[theorem]{Lemma}
\newtheorem{proposition}[theorem]{Proposition}
\newtheorem{corollary}[theorem]{Corollary}
\theoremstyle{definition}
\newtheorem{definition}[theorem]{Definition}
\theoremstyle{remark}
\newtheorem{remark}[theorem]{Remark}
\numberwithin{equation}{section}
\titleformat{\section}{\Large\bfseries\color{Ink}}{\thesection}{0.7em}{}
\titleformat{\subsection}{\large\bfseries\color{Ink}}{\thesubsection}{0.7em}{}
\titlespacing*{\section}{0pt}{2.4ex plus 1ex minus .2ex}{1.2ex plus .2ex}
\titlespacing*{\subsection}{0pt}{2.0ex plus .7ex minus .2ex}{.8ex plus .2ex}
\setlist{itemsep=.35em,topsep=.5em,leftmargin=*}
\setlength{\parskip}{.3em}
\setlength{\parindent}{1.25em}
\setlength{\emergencystretch}{1.5em}
\allowdisplaybreaks[1]
\raggedbottom
\pagestyle{fancy}
\fancyhf{}
\fancyhead[L]{\small\color{Muted}ProveIt research continuation}
\fancyhead[R]{\small\color{Muted}8 October 2026}
\fancyfoot[C]{\small\thepage}
\renewcommand{\headrulewidth}{.3pt}
\setcounter{tocdepth}{2}

\begin{document}
\begin{titlepage}
\thispagestyle{empty}
\setlength{\parindent}{0pt}
{\large\sffamily\color{Teal}PROVEIT \quad / \quad ANALYSIS \quad / \quad POLYLOGARITHMS\par}
\vspace{20mm}
{\huge\bfseries\color{Ink}
Real Zeros and Exact Relation Spaces\par}
\vspace{4mm}
{\LARGE\bfseries\color{Ink}
in the Polylogarithm--Stieltjes Programme\par}
\vspace{8mm}
{\large
Research continuation, asymptotic expansions,\\
and corrections to the ProveIt drafts\par}
\vspace{13mm}
{\normalsize
Prepared for Vladimir Reshetnikov\\
Research and computational assistance: OpenAI ChatGPT\\
8 October 2026\par}
\vspace{12mm}
{\color{Teal}\rule{\linewidth}{.8pt}\par}
\vspace{3mm}
\noindent\textbf{Abstract.}
We continue the polylogarithm and special-function programme in the
ProveIt draft corpus. A reciprocal-gamma Appell polynomial provides
an exact Laplace kernel for every positive parameter derivative of
the generalized Stieltjes constants. We prove that
$\gamma_n^{(k)}(a)$ has at most $n$ positive real zeros for every
$k\geq1$, that all $n$ simple zeros occur for sufficiently large $k$,
and that their locations admit complete asymptotic expansions.
For the first Laurent index, we prove uniqueness for every derivative
order and a uniform half-unit bracket. A conductor-character
argument proves the unrestricted formal rational-grid rank laws,
including arbitrary weights and free finite-jet modules over
coefficient algebras. We give the exact functional-equation bridge
at trivial Dirichlet-$L$ zeros and reconstruct the known all-even-weight
Gaussian Euler-sum reduction with attribution. The audit supplies
explicit corrections to convergence, depth, iterated-integral,
modified-polylogarithm, log-gamma, and Herglotz claims; exact
class-number-one certificates repair five cubic examples.
Reproducible exact and high-precision checks accompany the proofs.
Thirteen further research questions distinguish quantitative
extensions, arithmetic independence, and formalization.

\vfill
{\small\color{Muted}
Source snapshot:
\texttt{3a6d80ed618d839915deb0d19ab687f45e81d995}.\\
An unrefereed research draft with explicit proofs and computational
witnesses; historical priority is stated conservatively.\par}
\end{titlepage}
\pagenumbering{roman}
\begingroup
\small
\setlength{\parskip}{0pt}
\tableofcontents
\endgroup
\clearpage
\pagenumbering{arabic}

% ----- sections/introduction.tex -----
\section{Scope, results, and relation to the existing drafts}
\label{intro:section}

This article continues the research programme in
\path{Analysis/Polylogarithms/docs} of the ProveIt repository.
The inspected snapshot is commit
\texttt{3a6d80ed618d839915deb0d19ab687f45e81d995},
dated 8 October 2026 UTC.  It contains eight articles and thirty-one
reports, covering multiple polylogarithms, cyclotomic values, gamma and
polygamma reductions, Herglotz functions, and generalized Stieltjes
constants \cite{RepoSnapshot}.  The source inventory and exact Git blob
identifiers are recorded in \path{PROVENANCE.json}.

The strongest analytic development here concerns real zeros of the
Stieltjes parameter derivatives.  The existing derivative-tower article
\cite{DraftTower} expresses these functions through Hurwitz zeta jets.
We reorganize that identity as a Laplace transform whose sign-changing
polynomial is independent of the derivative order.  This makes a global
zero bound and a complete large-order description possible.
The strongest algebraic development is a proof of the unrestricted
rational-grid rank formulas that were stated or conjectured from finite
matrices in the two Stieltjes articles \cite{DraftTower,DraftLadder}.

\subsection{The main analytic statements}

Our normalization is
\begin{equation}\label{intro:stieltjes}
 \zeta(s,a)=\frac1{s-1}
       +\sum_{n=0}^{\infty}\frac{(-1)^n}{n!}\gamma_n(a)(s-1)^n,
 \qquad a>0,
\end{equation}
and $\gamma_n^{(k)}$ always means a derivative in $a$.
For fixed Laurent index $n$, define
\[
 P_n(x)=n![z^n]\frac{e^{xz}}{\Gamma(1+z)}.
\]
The results of Sections~\ref{zero:section}--\ref{zero:asymptotic-section}
are the following.

\begin{enumerate}
\item For every $k\geq1$, $\gamma_n^{(k)}$ has at most $n$ positive real
zeros, counted with multiplicity.  This is a global statement on the
whole half-line, not a count inside a numerically chosen window.
\item The polynomial $P_n$ has $n$ distinct real roots.  For every fixed
$n$, the zero bound is attained for all sufficiently large $k$; all
zeros are simple and correspond bijectively to those roots.
\item If $\rho$ is a root of $P_n$ and $t=e^\rho$, its associated zero is
\begin{equation}\label{intro:zero-result}
 a_{\rho,k}=(k+1)e^{-\rho}
   +\frac{e^{-\rho}}2\left(\frac{P_n''(\rho)}{P_n'(\rho)}-1\right)
   -\frac1{\exp(e^\rho)-1}+O_n(k^{-1}).
\end{equation}
We give an explicit $1/(k+1)$ correction, an $O_n(k^{-2})$ remainder
after including it, and a recurrence for the full asymptotic expansion.
\item At $n=1$ there is a unique positive zero for every $k\geq1$, with
the global bracket
\[
 e^{H_{k-1}}<a_k<e^{H_{k-1}}+\tfrac12.
\]
It increases strictly with $k$.  More generally, once all $n$ simple
zeros are present, successive derivative orders strictly interlace.
\end{enumerate}

The elementary mechanism is worth emphasizing.  A reciprocal-gamma
Appell polynomial determines the sign changes of the Laplace kernel.
Rolle's theorem bounds how many zeros its transform can have.  A gamma
probability density then concentrates at $t=(k+1)/a$, yielding all
nearby zeros and their expansions.  Combining the local concentration
argument with the global count excludes zeros on unexamined scales.

\subsection{The formal relation spaces}

For the rational grids, one must first specify which values are unknown
and which are prescribed on the right-hand sides.  Our endpoint symbol
represents the value at $a=1$, not the singular limit at zero.
We prove a weighted distribution theorem for arbitrary rational prime
weights, including zero and resonant weights.  The resulting formal
quotient has dimension $\varphi(q)$ before prescribing the endpoint and
$\varphi(q)-1$ afterwards.  The proof identifies each character block
with its primitive-conductor coordinate and does not divide by an
Euler factor that might vanish.

The same construction works over every commutative coefficient algebra
over a characteristic-zero splitting field. Applied to
$K[\eta]/(\eta^{N+1})$, it proves freeness of the entire finite jet
presentation, with dimension $(N+1)(\varphi(q)-1)$ over $K$ after all
endpoint layers are prescribed. Here $K$ contains the relevant character
values and prime logarithms. This statement concerns the specified
distribution rows alone; additional evaluated identities can reduce a
larger arithmetic presentation.

In particular, the positive Stieltjes derivative grids have formal
residual dimension $\varphi(q)-1$.  For negative zeta jets, after
including the specified reflection identities and endpoint, the
dimension is
\[
 r_{q,k}=\begin{cases}
 \varphi(q)/2-1,& k\ \text{odd},\\
 \varphi(q)/2,& k\ \text{even},
 \end{cases}\qquad(q\geq3).
\]
These statements hold for every positive $k$ and every denominator,
not merely the finite ranges in the drafts.  They are dimensions of
formal presentations.  Injectivity after evaluation as actual complex
numbers is an additional arithmetic question.

We also complete the derivative-tower article's description at trivial
zeros with the exact identity
\[
 \frac{L'(k+1,\chi)}{L(k+1,\chi)}
 +\frac12\overline{\frac{L''(-k,\chi)}{L'(-k,\chi)}}
 =\gamma+\log(2\pi/q)-H_k,
 \qquad\chi(-1)=(-1)^k.
\]
The factor $1/2$ comes from cancelling the gamma pole against the
simple $L$-zero before taking a logarithmic derivative.

\subsection{Classical results, new proofs, and historical priority}

Several central ingredients predate these drafts.  Coffey gives the
general parameter-derivative formula and proves the unique maximum of
$\gamma_1(a)$, the $n=k=1$ case of our zero discussion
\cite{Coffey2009}.  Kubert's universal ordinary distribution gives the
classical rank-$\varphi(q)$ context \cite{Kubert1979}.  The Gaussian
all-even-weight reduction, including the draft's $S_5$ and $S_7$
specializations, appears in Olaikhan's 2022 derivation
\cite{gau-olaikhan-2022}.  Section~\ref{gau:section} gives a complete
analytic reconstruction with an exact generating function and a
convergent resummation, with that attribution retained.

The present continuation establishes the real-zero and asymptotic
theorems just stated, closes the two unrestricted formal-rank gaps by an
explicit weighted proof, and supplies correction proofs and witnesses.
The literature search was targeted to these statements and their
ingredients; it was not an exhaustive priority review.  We therefore
claim the proofs and their application to the repository, without
asserting that every corollary is absent from all earlier literature.
The mathematical arguments were checked in separate derivations during
preparation.  This remains an unrefereed research draft, without a Lean
formalization.

\subsection{Audit policy and source scope}

The audit concerns the retrieved \LaTeX{} and text sources.  All eight
article sources were examined, together with the reports supporting the
claims discussed below.  The remaining reports were inventoried for
context; no assertion of exhaustive verification of every displayed
formula in all thirty-nine documents is made.
Section~\ref{audit:section} provides explicit locations, counterexamples,
or replacement arguments.  The accompanying correction register
distinguishes false formulas from missing hypotheses, incomplete proofs,
and unproved numerical-independence claims.

This distinction matters in this corpus.  An exact rational matrix rank
proves something about that matrix.  It does not prove that every
numerical relation has been included.  A negative integer-relation
search does not establish transcendence or independence.  Nor does a
positive dimension of an ambient motivic space identify the classes of
specific named periods.  The corrected theorems keep these three issues
separate while preserving the valid identities and computations.

\subsection{Notation}

We use $H_k=\sum_{j=1}^k j^{-1}$, $H_0=0$,
$\beta(s)=\sum_{r\geq0}(-1)^r(2r+1)^{-s}$, and $\varphi$ for Euler's
totient.  A prime on $L$ or $\zeta$ means differentiation in the complex
variable $s$ unless a parameter derivative is explicitly indicated.
For $\nu\geq2$, the Clausen convention is
\[
 \operatorname{Cl}_\nu(\theta)=
 \begin{cases}
  \Im\Li_\nu(e^{i\theta}),&\nu\ \text{even},\\
  \Re\Li_\nu(e^{i\theta}),&\nu\ \text{odd}.
 \end{cases}
\]
The Herglotz function in the audit is
$F(x)=\sum_{m\geq1}(\psi(mx)-\log(mx))/m$, unrelated to
$F_{n,k}$ except for the letter used in its source literature.
Each section restates the hypotheses needed for its own formulas.

% ----- sections/stieltjes_zeros.tex -----
\section{A polynomial kernel for the Stieltjes derivative tower}
\label{zero:section}

The rank calculation describes relations at rational arguments.  A different
question is how the functions themselves change sign as the positive real
argument varies.  Here the derivative order $k$ provides a concentration
parameter.  We keep the Laurent index $n$ fixed and let $k$ grow; this is
different from the frequently studied asymptotics of $\gamma_n(a)$ as
$n\to\infty$.

Define monic polynomials $P_n$ by
\begin{equation}\label{zero:appell}
 \frac{e^{xz}}{\Gamma(1+z)}
   =\sum_{n=0}^{\infty}P_n(x)\frac{z^n}{n!}.
\end{equation}
The generating function is entire in $z$.  With $y=x+\gamma$, its first
members are
\begin{align*}
 P_0(x)&=1, & P_1(x)&=y,\\
 P_2(x)&=y^2-\zeta(2),&
 P_3(x)&=y^3-3\zeta(2)y+2\zeta(3),\\
 P_4(x)&=y^4-6\zeta(2)y^2+8\zeta(3)y
                       +3\zeta(2)^2-6\zeta(4).
\end{align*}
They satisfy $P_n'=nP_{n-1}$ and the constructive recurrence
\begin{equation}\label{zero:recurrence}
 P_{n+1}(x)=(x+\gamma)P_n(x)
       +\sum_{r=1}^{n}(-1)^r\frac{n!}{(n-r)!}\zeta(r+1)P_{n-r}(x).
\end{equation}
This follows by differentiating \eqref{zero:appell} in $z$ and using the
Taylor series of $\log\Gamma(1+z)$ near zero.  These are the Appell
polynomials associated with the reciprocal gamma function; the use of
real-rooted Appell polynomials belongs to classical Laguerre--P\'olya
theory.  We supply the needed real-root and multiplicity argument below,
so no classification theorem is required.

\begin{theorem}[Exact Laplace kernel]\label{zero:kernel}
For $n\geq0$, $k\geq1$, and $a>0$, set
$F_{n,k}(a)=(-1)^{n+k}\gamma_n^{(k)}(a)$.  Then
\begin{equation}\label{zero:laplace}
 F_{n,k}(a)=\int_0^\infty
       \frac{t^k e^{-at}}{1-e^{-t}}P_n(\log t)\,dt.
\end{equation}
The integral and all its derivatives in $a$ converge locally uniformly
on $a>0$, and $F_{n,k}'=-F_{n,k+1}$.
\end{theorem}
\begin{proof}
Start from the classical Mellin integral for Hurwitz zeta
\cite{DLMF25}.  Differentiating $k$ times in the parameter gives, initially
for $\Re s>1$,
\[
 \partial_a^k\zeta(s,a)
   =\frac{(-1)^k}{\Gamma(s)}
      \int_0^\infty\frac{t^{s+k-1}e^{-at}}{1-e^{-t}}\,dt.
\]
For $k\geq1$ the right side is holomorphic in a neighborhood of $s=1$:
near $t=0$ its integrand is $O(t^{k-1+\Re(s-1)})$, and infinity is
controlled by $e^{-at}$.  The parameter derivative of the pole
$1/(s-1)$ is zero.  Put $s=1+z$, insert \eqref{zero:appell}, and
compare the coefficient of $z^n$ with
$(-1)^n\gamma_n^{(k)}(a)/n!$.  The bounds just stated, with finitely
many logarithmic powers or extra factors of $t$, justify every
coefficient extraction and parameter differentiation.
\end{proof}

This representation is compatible with Coffey's derivative formulas
\cite{Coffey2009} and the harmonic-number master identity in the draft.
Its useful feature here is that $P_n$ is independent of $k$: the entire
dependence on derivative order is the positive multiplier $t^k$.

\subsection{Why all kernel roots are real and simple}

\begin{lemma}\label{zero:root-preserver}
Let $p$ be a monic real polynomial with only real roots, and let $c>0$.
Then $p+cp'$ has only real roots.  A root of $p$ of multiplicity $m$
has multiplicity $m-1$ in $p+cp'$ if $m>1$ and is absent if $m=1$.
Every remaining root is simple.  Applying $n-1$ operators of the form
$1+cD$ to any real-rooted degree-$n$ polynomial therefore gives a
polynomial with $n$ distinct real roots.
\end{lemma}
\begin{proof}
If the distinct roots are $r_1<\cdots<r_\ell$, with multiplicities
$m_i$, then away from them
\[
 \frac{p'(x)}{p(x)}=\sum_{i=1}^{\ell}\frac{m_i}{x-r_i},
 \qquad
 \left(\frac{p'}p\right)'=-\sum_{i=1}^{\ell}
                                  \frac{m_i}{(x-r_i)^2}<0.
\]
The equation $p'/p=-1/c$ has exactly one simple solution in each
$(r_i,r_{i+1})$ and one in $(-\infty,r_1)$, and none in
$(r_\ell,\infty)$.  At each old root the multiplicity decreases
by one, as direct factorization shows.  This accounts for all $n$
roots.  Iteration reduces the largest possible multiplicity to one.
\end{proof}

\begin{theorem}[Simple kernel roots]\label{zero:roots}
For every $n\geq1$, $P_n$ has exactly $n$ distinct real roots.
The roots of $P_{n-1}$ strictly interlace those of $P_n$.
\end{theorem}
\begin{proof}
The Weierstrass product \cite{DLMF5} is
\[
 \frac1{\Gamma(1+z)}
   =e^{\gamma z}\prod_{j=1}^{\infty}(1+z/j)e^{-z/j}.
\]
Its finite products converge uniformly on compact sets.  At degree $n$
the corresponding Appell polynomial is
\[
 \prod_{j=1}^{M}(1+D/j)(x+\gamma-H_M)^n.
\]
Lemma~\ref{zero:root-preserver} makes it real-rooted, and convergence
of coefficients shows that $P_n$ is real-rooted and monic.  This limit
argument alone does not prove simplicity, since roots could coalesce.

To exclude coalescence, factor the first $n-1$ linear factors out of
the entire function, leaving
\[
 g_n(z)=e^{(\gamma-H_{n-1})z}
             \prod_{j=n}^{\infty}(1+z/j)e^{-z/j}.
\]
Its degree-$n$ Appell polynomial $Q_n(x)=n![z^n]e^{xz}g_n(z)$ is
again a coefficient limit of monic real-rooted polynomials.  Hence it
is real-rooted.  But now the exact finite identity
\[
 P_n=\prod_{j=1}^{n-1}(1+D/j)Q_n
\]
and Lemma~\ref{zero:root-preserver} give simplicity after taking
the limit.  The case $n=1$ is immediate.  Finally
$P_n'=nP_{n-1}$ and Rolle's theorem give strict interlacing.
\end{proof}

\section{A global bound on positive real zeros}
\label{zero:count-section}

We use a short form of the variation-diminishing principle for Laplace
transforms, proved here with Rolle's theorem.  A sign change means a
change between strict positive and negative signs; isolated zeros
without a sign change are not counted as sign changes.

\begin{lemma}[Laplace sign-variation bound]\label{zero:variation}
Let $f$ be a nonzero continuous real-valued function on $(0,\infty)$ with exactly
$m$ sign changes and with
$\int_0^\infty e^{-at}(1+t^r)|f(t)|\,dt<\infty$ for every $a>0$
and nonnegative integer $r$.  Assume $f$ is nonzero on a
set of positive measure.  Its Laplace transform
$L_f(a)=\int_0^\infty e^{-at}f(t)\,dt$ has at most $m$ zeros on
$(0,\infty)$, counted with analytic multiplicity.
\end{lemma}
\begin{proof}
The integrability hypotheses give a real-analytic Laplace transform
and justify differentiation under its integral.
For $m=0$ the transform has a fixed strict sign.  For $m>0$ choose
one sign-change point $t_0$.  Multiplication by $t_0-t$ removes
that sign change and creates no new one, so
$g(t)=(t_0-t)f(t)$ has $m-1$ sign changes.  Moreover,
\[
 \frac{d}{da}\bigl(e^{at_0}L_f(a)\bigr)=e^{at_0}L_g(a).
\]
If $L_f$ has $N$ zeros on a compact interval, counted with
multiplicity, the derivative on the left has at least $N-1$ zeros
there.  By induction $N-1\leq m-1$.  Applying this to any finite
collection of zeros proves the result on the entire half-line.
In particular, none of the transforms encountered in the induction
can vanish identically.
\end{proof}

\begin{theorem}[Uniform zero-count bound]\label{zero:at-most}
For every $n\geq0$ and every $k\geq1$, the real function
$a\mapsto\gamma_n^{(k)}(a)$ has at most $n$ positive real zeros,
counted with multiplicity.
\end{theorem}
\begin{proof}
The kernel in \eqref{zero:laplace}, before the factor $e^{-at}$,
is $t^kP_n(\log t)/(1-e^{-t})$.  Its positive factors do not change
sign, and Theorem~\ref{zero:roots} gives exactly $n$ sign changes.
All exponential moments exist by the estimates in the proof of
Theorem~\ref{zero:kernel}.  Apply Lemma~\ref{zero:variation}.
\end{proof}

The bound is independent of $k$, although the zero locations move to
infinity as $k$ increases.  We will prove that it is attained for
every sufficiently large $k$ at each fixed $n$.  There is no claim
here that it is attained at every small derivative order for $n>1$.

For reference, elementary endpoint estimates are
\begin{align}
 F_{n,k}(a)&\sim \Gamma(k+1)a^{-k-1}\bigl(\log(1/a)\bigr)^n,
                  &&a\downarrow0,\label{zero:left-end}\\
 F_{n,k}(a)&\sim (-1)^n\Gamma(k)a^{-k}(\log a)^n,
                  &&a\to\infty.\label{zero:right-end}
\end{align}
For \eqref{zero:left-end}, substitute $u=at$ and use
$(1-e^{-u/a})^{-1}\to1$ away from zero; the small-$u$ part is
controlled by $(1-e^{-v})^{-1}\leq1+1/v$.
For \eqref{zero:right-end}, the same substitution uses
$(1-e^{-u/a})^{-1}\sim a/u$ and the leading monic term of $P_n$.
Dominated convergence after division by the displayed logarithmic
power proves both statements; integrals of $u^{k-1}|\log u|^r e^{-u}$
are finite.  The interpretations for $n=0$ are immediate.

\section{Uniform expansions and the motion of every zero}
\label{zero:asymptotic-section}

Put $K=k+1$ and
\[
 h_n(t)=\frac{P_n(\log t)}{1-e^{-t}},\qquad
 \mathcal H_{n,K}(t)=\frac{(K/t)^K}{\Gamma(K)}F_{n,K-1}(K/t).
\]
If $U_K$ has gamma density
$K^K u^{K-1}e^{-Ku}/\Gamma(K)$ on $u>0$, then
\begin{equation}\label{zero:expectation}
 \mathcal H_{n,K}(t)=\mathbb E\,h_n(tU_K),
 \qquad \mathbb E U_K=1.
\end{equation}
This is an exact identity.  It identifies the useful scale $a=K/t$
and turns the asymptotic calculation into an expansion around $U_K=1$.

\begin{theorem}[All-order local expansion]\label{zero:all-orders}
Fix $n\geq0$, $M\geq0$, a compact interval $I\subset(0,\infty)$,
and a nonnegative integer $r$.  Uniformly for $t\in I$, together with
the first $r$ derivatives in $t$,
\begin{equation}\label{zero:expansion}
 \mathcal H_{n,K}(t)=\sum_{j=0}^{M}\frac{\mathcal D_jh_n(t)}{K^j}
                         +O_{n,M,r,I}(K^{-M-1}).
\end{equation}
Here $\mathcal D_0h=h$, and for $j\geq1$ the finite formula is
\begin{equation}\label{zero:D-general}
 \mathcal D_jh(t)=
 \sum_{\substack{m_2,m_3,\ldots\geq0\\
                 \sum_{\ell\geq2}(\ell-1)m_\ell=j}}
 \frac{t^R h^{(R)}(t)}{
             \prod_{\ell\geq2}m_\ell!\,\ell^{m_\ell}},
 \qquad R=\sum_{\ell\geq2}\ell m_\ell.
\end{equation}
In particular,
\begin{align}
 \mathcal D_1h&=\tfrac12t^2h'',\label{zero:D1}\\
 \mathcal D_2h&=\tfrac13t^3h'''+\tfrac18t^4h^{(4)},\label{zero:D2}\\
 \mathcal D_3h&=\tfrac14t^4h^{(4)}+\tfrac16t^5h^{(5)}
                                      +\tfrac1{48}t^6h^{(6)}.
\end{align}
The constants in the remainder can be bounded using suprema of finitely
many derivatives of $h_n$ on an enlarged compact interval and explicit
gamma moments and tail bounds.  No uniformity in an increasing $n$
is asserted.
\end{theorem}
\begin{proof}
The cumulants of $U_K-1$ are zero in degree one and
$(\ell-1)!K^{1-\ell}$ in degree $\ell\geq2$.  Expanding its moment
generating function therefore gives the exact central-moment identity
\begin{equation}\label{zero:moments}
 \frac{\mathbb E(U_K-1)^R}{R!}
 =\sum_{\substack{\sum_{\ell\geq2}\ell m_\ell=R}}
       \prod_{\ell\geq2}
       \frac{K^{-(\ell-1)m_\ell}}{m_\ell!\,\ell^{m_\ell}}.
\end{equation}
For $R=0$ the value is one, and for $R=1$ it is zero.
When $R\geq2$, the smallest possible power of $K^{-1}$ is
$\lceil R/2\rceil$.

Expand $h_n(tu)$ about $u=1$ to degree $2M+1$, on
$1/2\leq u\leq2$.  Its remainder is bounded uniformly by a constant
times $|u-1|^{2M+2}$.  The expectation of this even power is
$O(K^{-M-1})$ by \eqref{zero:moments}.  The coefficient of $K^{-j}$
in the expected Taylor polynomial is precisely
\eqref{zero:D-general}; terms with $j>M$ are $O(K^{-M-1})$.

For completeness, the discarded tails are exponentially small, even
though $h_n$ is unbounded at zero.  Chernoff's bound for the gamma
moment generating function gives, for fixed $v\ne1$ on its relevant
side of one,
\[
 \Pr(U_K\geq v)\ \text{or}\ \Pr(U_K\leq v)
       \leq \exp\{-K(v-1-\log v)\}.
\]
For each fixed derivative order, $u^r h_n^{(r)}(tu)$ is bounded
uniformly on $t\in I$ by $C(u^{-B}+u^B)$ for some fixed $B$.
The gamma moments
$\mathbb E U_K^v=K^{-v}\Gamma(K+v)/\Gamma(K)$ are uniformly bounded
for fixed positive or negative $v$ once $K$ is sufficiently large.
Cauchy--Schwarz therefore bounds the tail expectations by
$Ce^{-cK}$.  The same reasoning replaces truncated central moments
by their full values.  Finally differentiate
\eqref{zero:expectation} in $t$, obtaining
$\mathbb E[U_K^r h_n^{(r)}(tU_K)]$, and repeat the argument.
This proves the stated $C^r$ uniformity and the expansion.
\end{proof}

\begin{theorem}[Saturation and asymptotic locations]\label{zero:locations}
Fix $n\geq1$, and let $\rho$ run over the distinct real roots of $P_n$.
For all sufficiently large $k$, $\gamma_n^{(k)}$ has exactly $n$
positive real zeros, all simple.  One zero is associated with each
$\rho$.  With $t_\rho=e^\rho$, it has the expansion
\begin{equation}\label{zero:location-first}
 a_{\rho,k}=\frac{k+1}{t_\rho}
       +\frac{1}{2t_\rho}
           \left(\frac{P_n''(\rho)}{P_n'(\rho)}-1\right)
       -\frac1{e^{t_\rho}-1}+O_n(k^{-1}).
\end{equation}
Each zero has a full expansion in descending integer powers of $k+1$.
In particular $a_{\rho,k}/k\to e^{-\rho}$.
\end{theorem}
\begin{proof}
The zeros of $h_n$ are the simple points $t_\rho$.  By the uniform
$C^1$ version of Theorem~\ref{zero:all-orders}, each has a disjoint
neighborhood containing a unique simple zero $t_{\rho,K}$ of
$\mathcal H_{n,K}$ for all sufficiently large $K$.  Thus there are
at least $n$ positive zeros.  Theorem~\ref{zero:at-most} excludes
every additional zero, including any outside the compact intervals
used in this local argument.

Write $t=t_\rho$, $h=h_n$, and
$A(t)=t^2h''(t)/2$.  Expansion of the zero equation gives
\[
 t_{\rho,K}=t-\frac{A(t)}{K h'(t)}+O(K^{-2}),
 \qquad
 a_{\rho,k}=\frac K{t_{\rho,K}}
              =\frac Kt+\frac{h''(t)}{2h'(t)}+O(K^{-1}).
\]
At a root of $P_n(\log t)$, direct differentiation gives
\[
 \frac{h''(t)}{h'(t)}
  =\frac1t\left(\frac{P_n''(\rho)}{P_n'(\rho)}-1\right)
                           -\frac2{e^t-1},
\]
which proves \eqref{zero:location-first}.  At every subsequent order
the new coefficient of $t_{\rho,K}$ is multiplied by the same nonzero
number $h'(t)$; the all-order expansion and a Taylor remainder therefore
determine it recursively.  Applying the mean value theorem to the
residual of a truncated zero expansion proves the stated remainder at
each finite order.
\end{proof}

The distinction between local concentration and global zero counting is
essential: concentration alone finds $n$ nearby zeros but cannot rule
out additional zeros on scales escaping the chosen compact interval.
The variation bound supplies precisely that missing global step.

\paragraph{An explicit next correction.}
Let $b_j=h_n^{(j)}(t_\rho)/h_n'(t_\rho)$ for $j=2,3,4$ and set
\begin{equation}\label{zero:next-coefficient}
 c_\rho=t_\rho\left(\frac{b_3}{3}-\frac{b_2^2}{4}\right)
      +\frac{t_\rho^2}{8}
                    \left(b_4-2b_2b_3+b_2^3\right).
\end{equation}
Then the more accurate formula is
\begin{equation}\label{zero:location-second}
 a_{\rho,k}=\frac{k+1}{t_\rho}+\frac{b_2}{2}
                         +\frac{c_\rho}{k+1}+O_n(k^{-2}).
\end{equation}
Indeed, with $B=\mathcal D_2h$, the first two coefficients of
$t_{\rho,K}=t+u_1/K+u_2/K^2+\cdots$ satisfy
$u_1=-A/h'$ and
$u_2=-(h''u_1^2/2+A'u_1+B)/h'$.  Expanding $K/t_{\rho,K}$
and collecting terms gives \eqref{zero:next-coefficient}.

\begin{corollary}[Persistence and interlacing after saturation]
\label{zero:k-interlace}
Fix $n\geq1$. If $\gamma_n^{(k)}$ has $n$ distinct positive
zeros at some $k\geq1$, then the same is true at order $k+1$.
Writing the zeros in increasing order, one has
\[
 a_{j,k}<a_{j,k+1}<a_{j+1,k}\quad(1\leq j<n),
 \qquad a_{n,k}<a_{n,k+1}.
\]
In particular, these conclusions hold at every sufficiently large $k$.
\end{corollary}
\begin{proof}
Rolle's theorem and $F_{n,k}'=-F_{n,k+1}$ give a zero between each
consecutive pair of zeros of $F_{n,k}$.  Beyond its last zero,
$F_{n,k}$ has a fixed nonzero sign and tends to zero by
\eqref{zero:right-end}; it therefore has an interior extremum there.
This gives $n$ distinct zeros of $F_{n,k+1}$.  Its zero bound
excludes all others, including coincidences with a simple zero of
$F_{n,k}$, and establishes the claimed ordering.
\end{proof}

\subsection{The first Laurent index: a global half-unit bracket}

For $n=1$, the kernel has its sole sign change at $t_0=e^{-\gamma}$.
Here the conclusions can be made global in the derivative order.
Coffey already proved that $\gamma_1$ has a unique maximum at
$a\approx1.39112$ \cite[Proposition 5]{Coffey2009}; this is the
$k=1$ uniqueness case below.  The present statement extends the
description to every positive derivative order and supplies the
uniform bracket and moving-zero expansion.

\begin{theorem}[Unique zero and an explicit bracket]\label{zero:first}
For every integer $k\geq1$, $\gamma_1^{(k)}(a)$ has a unique positive
zero $a_k$, which is simple.  With $H_0=0$,
\begin{equation}\label{zero:half-unit}
 e^{H_{k-1}}<a_k<e^{H_{k-1}}+\tfrac12.
\end{equation}
The zeros are strictly increasing in $k$.  Their asymptotic expansion
starts with
\begin{equation}\label{zero:first-asymptotic}
 a_k=e^\gamma\left(k+\tfrac12\right)
                     -\frac1{e^{e^{-\gamma}}-1}+O(k^{-1}).
\end{equation}
Formula~\eqref{zero:location-second} supplies the next coefficient
and an $O(k^{-2})$ remainder.
\end{theorem}
\begin{proof}
The endpoint signs in \eqref{zero:left-end}--\eqref{zero:right-end}
and the bound of one zero prove existence, uniqueness, and simplicity.
To obtain the bracket, normalize the positive density
$t^k e^{-at}/(1-e^{-t})$.  The sign of $F_{1,k}(a)$ is the sign of
$\mathbb E(\log T+\gamma)$ under this density.

Relative to a gamma density of shape $k$ and rate $a$, the tilting
factor is
$r(t)=t/(1-e^{-t})$, which is strictly increasing: its derivative has
numerator $1-e^{-t}(1+t)>0$.  A strictly increasing tilt strictly
increases the expectation of the strictly increasing function
$\log t$.  This elementary fact follows from
\[
 2\operatorname{Cov}(u(X),v(X))
   =\mathbb E[(u(X)-u(Y))(v(X)-v(Y))]>0
\]
for independent $X,Y$ with the same continuous distribution and
strictly increasing $u,v$.  Consequently
\[
 \mathbb E\log T>\psi(k)-\log a
                 =H_{k-1}-\gamma-\log a.
\]
At the zero this gives $a_k>e^{H_{k-1}}\geq1$.

For $a>1/2$, compare instead with a gamma density of shape $k$
and rate $a-1/2$.  The tilting factor is
$s(t)=t/(2\sinh(t/2))$, which is strictly decreasing because
$(t/2)\cosh(t/2)>\sinh(t/2)$.  The covariance inequality reverses:
\[
 \mathbb E\log T<\psi(k)-\log(a-1/2).
\]
Applying it at $a_k>1$ gives $a_k<e^{H_{k-1}}+1/2$.

Finally, at a zero the signed kernel $f(t)$ in
\eqref{zero:laplace} satisfies $\int e^{-a_kt}f(t)\,dt=0$ and
$(t-t_0)f(t)>0$ except at $t_0$.  Thus
\[
 F_{1,k+1}(a_k)=\int_0^\infty (t-t_0)e^{-a_kt}f(t)\,dt>0.
\]
Since $F_{1,k+1}$ crosses from positive to negative at its unique
zero, $a_{k+1}>a_k$.  The root of $P_1$ is $-\gamma$, so
\eqref{zero:first-asymptotic} is Theorem~\ref{zero:locations}.
\end{proof}

The interval in \eqref{zero:half-unit} is an analytic certificate for
bracketing a unique root, valid even at $k=1$.  Its width is uniformly
$1/2$; no assertion that this is the best possible width is made.
The accompanying computation uses this bracket rather than guessing
an initial zero from a graph.

\begin{figure}[tbp]
\centering
\includegraphics[width=0.98\linewidth]{figures/stieltjes_profiles.pdf}
\caption{Normalized derivative profiles
$(1-e^{-e^x})\mathcal H_{n,K}(e^x)$ approach $P_n(x)$ as $K=k+1$
increases.  The plots illustrate $n=1,2,3$.  Quadrature produces the
curves; Theorems~\ref{zero:at-most} and \ref{zero:locations}, rather
than the plots, prove the zero counts and asymptotic completeness.}
\label{zero:profile-figure}
\end{figure}

% ----- sections/distribution_ranks.tex -----
\section{A complete proof of the rational-grid rank laws}
\label{rank:section}

The derivative-tower draft states an all-denominator rank law but supports it
only by finite exact computations.  The antiderivative draft states the
corresponding finite-range result carefully and leaves its extension to all
denominators as a conjecture.  Both questions admit the same elementary
character-theoretic solution.  The result concerns the rank of an explicitly
specified \emph{formal presentation}; it does not establish linear
independence of numerical zeta or Stieltjes values.

The relevant classical context is the universal ordinary distribution.
Kubert proved, in particular, that its level-$q$ module is free of rank
$\varphi(q)$ and becomes the regular unit-group representation over
$\mathbb Q$ \cite{Kubert1979}.  The proof below treats arbitrary rational
prime weights, states the endpoint convention explicitly, and obtains the
two rank formulas needed by the drafts.  We make no claim that the
underlying universal-distribution principle is new.

\subsection{The presentation, including the endpoint}

Fix $q\geq1$ and set
\[
 X_q=\tfrac1q\mathbb Z/\mathbb Z,
 \qquad V_q=\bigoplus_{x\in X_q}\mathbb Q e_x.
\]
When symbols are evaluated by a special function, use the representative
$\langle x\rangle_+\in(0,1]$.  In particular $e_0$ represents its value at
$1$, not its singular limit at $0$.  Choose a rational weight $w_p$ for each
prime $p\mid q$ and put
\[
 w(d)=\prod_{p\mid d}w_p^{v_p(d)}\qquad(d\mid q),
 \qquad w(1)=1.
\]
For every $d\mid q$ and $x\in X_{q/d}$ define a distribution row
\begin{equation}\label{rank:row}
 D_{d,x}=\sum_{\substack{y\in X_q\\dy=x}}e_y-w(d)e_x.
\end{equation}
There are exactly $d$ preimages in this sum.  Let $R_q(w)$ be the span of
these rows and let $U_q(w)=V_q/R_q(w)$.  The numerical application uses
$w_p=p^s$, with $s=k+1$ for positive parameter derivatives and $s=-k$ for
negative zeta jets.  Thus all the coefficient matrices are rational.

There are two distinct conventions worth recording.  The presentation
$U_q(w)$ retains the endpoint as a free symbol.  If the endpoint is placed
in a prescribed right-hand-side inventory, the homogeneous coefficient
space is
\begin{equation}\label{rank:endpoint}
 U_q^0(w)=V_q/(R_q(w)+\mathbb Q e_0).
\end{equation}
This is equivalent to moving every occurrence of the endpoint to the
right-hand side and using only the $q-1$ columns $e_{r/q}$,
$1\leq r<q$.

\begin{lemma}[Prime rows suffice]\label{rank:prime}
The rows $D_{p,x}$, for $p\mid q$ prime and $x\in X_{q/p}$, generate
$R_q(w)$.
\end{lemma}
\begin{proof}
For $ab\mid q$ and $x\in X_{q/(ab)}$, partition the $ab$ preimages by
their images under multiplication by $a$.  This gives the exact identity
\[
 D_{ab,x}=\sum_{\substack{y\in X_{q/a}\\by=x}}D_{a,y}
             +w(a)D_{b,x}.
\]
The factorization $w(ab)=w(a)w(b)$ holds even when $a$ and $b$ are not
coprime.  Repeated factorization proves the assertion, including all
prime-power rows.
\end{proof}

\subsection{The weighted character theorem}

Write $G_q=(\mathbb Z/q\mathbb Z)^\times$, acting by $u e_x=e_{ux}$.
For $q=1$ this is the trivial group.  If $\chi$ is a character of $G_q$,
write $f=f_\chi\mid q$ for its conductor and $\chi^*$ for its primitive
character.  The principal character has conductor $1$ and
$\chi^*(p)=1$.  For every $m$ with $f\mid m\mid q$, put
\begin{equation}\label{rank:E}
 E_{m,\chi}=
 \sum_{a\in(\mathbb Z/m\mathbb Z)^\times}
       \chi^*(a)e_{a/m}.
\end{equation}
At $m=1$ this is $e_0$.  These vectors are used after extending scalars to
$\mathbb C$; they need not individually have rational coefficients.

\begin{theorem}[Weighted distribution module]\label{rank:main}
For every $q\geq1$ and every choice of rational weights $w_p$, the
complexification of $U_q(w)$ has the basis
\[
 \bigl\{[E_{f_\chi,\chi}]:\chi\in\widehat G_q\bigr\}.
\]
The $\chi$-labelled basis vector has $G_q$-character $\chi^{-1}$.
Consequently
\begin{equation}\label{rank:dimensions}
 \dim_{\mathbb Q} U_q(w)=\varphi(q),
 \qquad
 \dim_{\mathbb Q} U_q^0(w)=\varphi(q)-1.
\end{equation}
In particular the first dimension is independent of all prime weights.
\end{theorem}

\begin{proof}
Decompose $V_q\otimes\mathbb C$ by the exact order $m\mid q$ of its
symbols.  The summand of order $m$ is the regular representation of
$G_m$, inflated to $G_q$.  The reduction map $G_q\to G_m$ is surjective
by the Chinese remainder theorem.  Hence a character $\chi$ of $G_q$
occurs in that summand if and only if $f_\chi\mid m$, and its
$\chi^{-1}$-eigenline is spanned by $E_{m,\chi}$.  This identifies every
character block before imposing relations.

By Lemma~\ref{rank:prime}, it is enough to sum prime distribution rows
against $\chi^*(a)$ over symbols of each exact order $m$.  Suppose
$pm\mid q$ and $f\mid m$.  If $p\mid m$, all $p$ lifts of a unit
$a/m$ have exact denominator $pm$.  Each unit modulo $pm$ occurs once,
with its correct character coefficient.  The resulting row is
\begin{equation}\label{rank:edge-ramified}
 E_{pm,\chi}-w_p E_{m,\chi}.
\end{equation}
If $p\nmid m$, exactly one lift has denominator $m$: writing that lift
as $b/m$, one has $pb\equiv a\pmod m$.  Its coefficient is
$\chi^*(a)=\chi^*(p)\chi^*(b)$.  The other lifts run through the units
modulo $pm$.  In this case the row is
\begin{equation}\label{rank:edge-unramified}
 E_{pm,\chi}-(w_p-\chi^*(p))E_{m,\chi}.
\end{equation}
The source space for rows of exact order $m$ has no $\chi$-component
when $f\nmid m$.  Since the distribution map is $G_q$-equivariant,
there are no further rows in the $\chi$-block.  This also explains why
a newly appearing primitive character is not removed when its conductor
is first reached.

For $f\mid m\mid q$ define
\begin{equation}\label{rank:A}
 A_{m,\chi}=
 \prod_{p\mid f}w_p^{v_p(m)-v_p(f)}
 \prod_{\substack{p\mid m\\p\nmid f}}
       (w_p-\chi^*(p))w_p^{v_p(m)-1}.
\end{equation}
Empty products and zeroth powers have value $1$, so $A_{f,\chi}=1$,
also if some weight is zero.  Multiplying the edge coefficients along
any divisor chain from $f$ to $m$ gives exactly $A_{m,\chi}$:
different prime steps commute.  Thus every class in this block satisfies
\begin{equation}\label{rank:transport}
 [E_{m,\chi}]=A_{m,\chi}[E_{f,\chi}],
\end{equation}
and the block has dimension at most one.

For the opposite inequality, define on the block the linear functional
$E_{m,\chi}\mapsto A_{m,\chi}$.  Formula~\eqref{rank:A} shows directly
that it vanishes on every edge row
\eqref{rank:edge-ramified}--\eqref{rank:edge-unramified}; its value on
$E_{f,\chi}$ is $1$.  The block therefore has dimension exactly one.
Summing over all $\varphi(q)$ characters proves the first formula in
\eqref{rank:dimensions}.  Rank does not change under the scalar extension
$\mathbb Q\subset\mathbb C$.  Finally $e_0=E_{1,\chi_0}$ is the nonzero
generator of the principal block, so adjoining $e_0$ to the relation
space removes exactly that block and proves the second formula.
\end{proof}

\begin{remark}[Resonance affects a choice of basis, not the rank]
\label{rank:resonance}
For $w_p=p^s$ with real $s\ne0$, every factor in
\eqref{rank:A} is nonzero, since $|\chi^*(p)|=1$ when $p\nmid f$.
For nonrational real weights the same proof is read over $\mathbb R$
instead of $\mathbb Q$, followed by extension to $\mathbb C$.
The classes of the $\varphi(q)$ unit symbols $e_{a/q}$ then form a basis
of $U_q(w)$: their Fourier coordinates are the nonzero multiples
$E_{q,\chi}=A_{q,\chi}E_{f,\chi}$.  At resonant weights this particular
basis can fail.  For example, if $w_p=1$ and $q=p$ is prime,
$E_{p,\chi_0}=0$ whereas $e_0$ still spans the principal block.
The conductor basis in Theorem~\ref{rank:main} remains valid, so there
is no exceptional rank.  No division by $w_p-\chi^*(p)$ occurs in the
proof.
\end{remark}

\subsection{The full finite-jet module is free}

The same argument controls all jet layers at once.  This gives more than
adding the dimensions of successive layers: it proves a free presentation
before the nilpotent jet parameter is specialized.

\begin{theorem}[Freeness over a coefficient algebra]\label{rank:algebra}
Let $K$ be a field of characteristic zero containing the values of all
characters of $G_q$, and let $A$ be any commutative $K$-algebra with
identity.  Define $V_q(A)$ and its distribution rows as in
\eqref{rank:row}, now with arbitrary weights $w_p\in A$ and with
$A$-linear spans.  Then
\[
 V_q(A)/R_q(w;A)\simeq A^{\varphi(q)},\qquad
 V_q(A)/(R_q(w;A)+Ae_0)\simeq A^{\varphi(q)-1}.
\]
The conductor vectors $E_{f_\chi,\chi}$ give the first basis; omitting
the principal vector gives the second.  These statements remain valid
when $A$ has zero divisors or nilpotents and when edge coefficients vanish.
\end{theorem}
\begin{proof}
The Fourier changes of basis over each $G_m$ and the character
projections divide only by $|G_m|$, which is invertible in $K$ and
therefore in $A$.  The decomposition into character blocks thus works
over $A$ exactly as it does over $\mathbb C$.
Within a block, the edge equations
\eqref{rank:edge-ramified}--\eqref{rank:edge-unramified} express every
generator as $A_{m,\chi}$ times the conductor generator.  They give a
surjection $A\to U_\chi$, $1\mapsto[E_{f_\chi,\chi}]$.
The map $E_{m,\chi}\mapsto A_{m,\chi}$ descends to $U_\chi\to A$
and is its inverse, since it sends the conductor generator to $1$.
Hence each block is free of rank one, without any use of division by
a weight or an Euler factor.  Removing the principal block gives the
endpoint assertion.
\end{proof}

\begin{corollary}[Finite Hurwitz and Stieltjes jets]\label{rank:full-jets}
Fix an integer $s\ne1$ and $N\geq0$, and let $K\subseteq\mathbb C$
contain all character values modulo $q$ and all $\log p$, $p\mid q$.
For $0\leq j\leq N$ use the formal symbols
$Z_j(x)=\zeta^{(j)}(s,\langle x\rangle_+)/j!$; the equality here
specifies their intended evaluation, not an arithmetic-independence
assumption.  Impose every coefficient of Hurwitz distribution through
order $N$, and prescribe all $N+1$ endpoint symbols $Z_j(0)$.
The resulting formal $K$-vector space has dimension
\begin{equation}\label{rank:full-jet-dimension}
 (N+1)(\varphi(q)-1).
\end{equation}
Before prescribing the endpoint its dimension is $(N+1)\varphi(q)$.
For fixed $k\geq1$, the same result holds for the complete finite
Stieltjes-derivative array
$\{\gamma_j^{(k)}(r/q):0\leq j\leq N,\ 1\leq r<q\}$,
with all its coefficient distribution rows and all endpoint layers
prescribed.
\end{corollary}
\begin{proof}
Use the truncated polynomial algebra $A=K[\eta]/(\eta^{N+1})$ and
the weights
\[
 w_p=p^s\sum_{\ell=0}^N\frac{(\log p)^\ell}{\ell!}\eta^\ell.
\]
They give $w(d)=d^s\exp(\eta\log d)$ in this algebra.  The coefficient
row at order $j$ is
\begin{equation}\label{rank:coefficient-jet-row}
 \sum_{dy=x}Z_j(y)
 -d^s\sum_{\ell=0}^j
       \frac{(\log d)^\ell}{\ell!}Z_{j-\ell}(x).
\end{equation}
Identify $Z_j(x)$ with $\eta^{N-j}e_x$ in the underlying $K$-vector
space of $V_q(A)$.  This reversed grading is necessary: multiplication
by $\eta$ lowers jet order.  Under this identification,
\eqref{rank:coefficient-jet-row} is exactly $\eta^{N-j}D_{d,x}$,
and prescribing all endpoint layers is exactly quotienting by $Ae_0$.
Theorem~\ref{rank:algebra} proves freeness over $A$; multiplying its
rank by $\dim_K A=N+1$ proves \eqref{rank:full-jet-dimension}.

Finally the coefficient identity
\[
 \sum_{j=0}^N\frac{(-1)^j}{j!}\gamma_j^{(k)}(a)\eta^j
 =(-1)^k(1+\eta)_k
       \sum_{j=0}^N\frac{\zeta^{(j)}(k+1,a)}{j!}\eta^j
       \pmod{\eta^{N+1}}
\]
is an invertible triangular change of jet coordinates: its scalar
factor has nonzero constant coefficient $(-1)^k k!$.  It preserves
the coefficient distribution presentation, with weights
$d^{k+1}\exp(\eta\log d)$, and proves the Stieltjes assertion.
\end{proof}

In particular, passing from an individual highest layer to the entire
finite array introduces no additional formal rank obstruction.  The
freeness result concerns these displayed relations over the specified
coefficient field; it still gives no numerical-independence theorem.

\subsection{Reflection and the negative-jet conjecture}

For $\epsilon\in\{1,-1\}$, let $S_q^\epsilon$ be generated by
$e_x+\epsilon e_{-x}$ for $x\in X_q$.  On the $\chi$-block, reflection
$x\mapsto-x$ acts by $\chi(-1)$, so these rows remove exactly the
characters satisfying $\chi(-1)=\epsilon$.

\begin{corollary}[Distribution and reflection rank]\label{rank:reflection}
For every $q\geq3$, every rational choice of weights, and
$\epsilon\in\{1,-1\}$,
\begin{equation}\label{rank:reflection-dimension}
 \dim_{\mathbb Q}
 \frac{V_q}{R_q(w)+S_q^\epsilon+\mathbb Q e_0}
 =\frac{\varphi(q)}2-\frac{1-\epsilon}{2}.
\end{equation}
For $q=1,2$ this dimension is $0$ for either sign.
\end{corollary}
\begin{proof}
The surviving characters have parity $-\epsilon$.  For $q\geq3$,
$-1$ is a nonidentity element of order two in $G_q$, so half of its
characters have each parity.  The principal character survives the
reflection quotient exactly when $\epsilon=-1$, and is then removed
by $e_0$.  For $q=1,2$, $G_q$ has only the principal character, which
the endpoint relation removes in either case.
\end{proof}

Here is the precise application to the drafts.  Put
$J_k(x)=\zeta'(-k,\langle x\rangle_+)$ for $k\geq1$.  Differentiating
Hurwitz multiplication gives
\begin{equation}\label{rank:jet-distribution}
 \sum_{j=0}^{d-1}\zeta'\left(-k,\frac{a+j}{d}\right)
 -d^{-k}\zeta'(-k,a)
 =d^{-k}\log d\,\zeta(-k,a).
\end{equation}
The reflection row used in the antiderivative draft has left-hand side
$J_k(x)+(-1)^kJ_k(-x)$.  More explicitly, for $0<x<1$ and
$C_\nu(x)+iS_\nu(x)=\operatorname{Li}_{\nu}(e^{2\pi i x})$, its known
right-hand side is
\begin{equation}\label{rank:jet-reflection}
 J_k(x)+(-1)^kJ_k(-x)
 =\begin{cases}
  \displaystyle\frac{(-1)^{k/2}k!}{(2\pi)^k}C_{k+1}(x),&k\text{ even},\\[6pt]
  \displaystyle\frac{(-1)^{(k-1)/2}k!}{(2\pi)^k}S_{k+1}(x),&k\text{ odd}.
 \end{cases}
\end{equation}
Indeed, differentiate the standard Hurwitz Fourier formula at $s=-k$;
the corresponding cosine or sine factor has a simple zero, so its
derivative is the sole surviving factor.  The endpoint value
$J_k(0)=\zeta'(-k)$ is placed in the inventory separately.

Take as the inventory the endpoint $\zeta'(-k)$, all Bernoulli-value
terms $d^{-k}\log d\,\zeta(-k,a)$ in
\eqref{rank:jet-distribution}, and all right-hand sides of
\eqref{rank:jet-reflection}.  Then
Corollary~\ref{rank:reflection}, with $s=-k$ and $\epsilon=(-1)^k$,
proves for \emph{all} $q\geq3$ and $k\geq1$ that the formal residual
dimension is
\begin{equation}\label{rank:negative-result}
 r_{q,k}=\begin{cases}
  \varphi(q)/2-1,&k\text{ odd},\\
  \varphi(q)/2,&k\text{ even}.
 \end{cases}
\end{equation}
This proves the unrestricted version of the antiderivative draft's
Proposition \texttt{prop:jetrank}, formerly checked only for
$3\leq q\leq30$ and $1\leq k\leq4$.  Equivalently, the remaining
complex character lines are indexed by primitive characters of
conductors $f\mid q$, $f>1$, with $\chi(-1)=(-1)^{k+1}$.

For the positive parameter tower, let the inventory contain
$\gamma_n^{(k)}(1)$ and every lower-$s$-derivative-layer term occurring
in the differentiated multiplication identity.  At a fixed highest
layer its homogeneous coefficient is $d^{k+1}$, independent of $n$.
In particular the $n=1$ row is
\[
 \sum_{j=0}^{d-1}\gamma_1^{(k)}\left(\frac{a+j}{d}\right)
 -d^{k+1}\gamma_1^{(k)}(a)
 =d^{k+1}\log d\,\psi^{(k)}(a).
\]
Theorem~\ref{rank:main} therefore proves the formal residual dimension
$\varphi(q)-1$ for every $q\geq1$ and $k\geq1$.  For $n>1$ this is
an associated-layer statement: lower derivatives are prescribed data,
not additional unknown columns.

\paragraph{What the theorem does and does not certify.}
The background inventory is part of the statement.  Several of its
terms, including rational-argument polygamma and Clausen values, are
not known to be elementary.  One may define an evaluation map from the
formal quotient to the span of actual constants modulo the inventory;
the theorem determines the dimension of its source.  Its image can
have smaller dimension, and proving injectivity would be a separate
arithmetic-independence problem.  Consequently ``one surviving formal
character line'' cannot be replaced by ``one numerically independent
new constant.''  The same caution applies to unsuccessful PSLQ searches.

\paragraph{Exact reproducibility.}
The accompanying \texttt{code/verify\_ranks.py} forms the actual
$q$-column rational matrices, including prime and all-divisor rows,
endpoint rows, and both reflection signs.  It checks their ranks by
exact fraction elimination and records its parameters and results in
\texttt{data/rank\_checks.json}.  It also checks the resonant weight
$s=0$ and a non-power weight system.  These computations are regression
checks on the presentation and its boundary conventions; the proof for
unbounded $q$ and $k$ is Theorem~\ref{rank:main}.
The companion \texttt{code/verify\_jet\_ranks.py} checks the expanded
finite-jet coefficient matrices with rational exponential-weight
parameters and with nilpotent weights; its output is
\texttt{data/jet\_rank\_checks.json}.

% ----- sections/bridge_completion.tex -----
\section{The exact bridge at a trivial zero}
\label{bridge:section}

The derivative-tower draft states the nonvanishing logarithmic bridge
correctly, but describes its trivial-zero analogue only as a ratio of
leading nonzero derivatives.  The precise coefficient matters: the
second-to-first derivative ratio appears with a factor of $1/2$.
The following statement supplies the omitted formula and proof.

\begin{theorem}[Regularized harmonic-number bridge]\label{bridge:trivial}
Let $\chi$ be a primitive Dirichlet character of conductor $q$, including
the conductor-$1$ character, and let $k\geq1$.  If
$\chi(-1)=(-1)^k$, then $L(s,\chi)$ has a simple zero at $s=-k$ and
\begin{equation}\label{bridge:formula}
 \frac{L'(k+1,\chi)}{L(k+1,\chi)}
 +\frac12\overline{\left(
       \frac{L''(-k,\chi)}{L'(-k,\chi)}\right)}
 =\gamma+\log\frac{2\pi}{q}-H_k.
\end{equation}
If instead $\chi(-1)=(-1)^{k+1}$, then $L(-k,\chi)\ne0$ and the
corresponding identity is
\begin{equation}\label{bridge:nontrivial}
 \frac{L'(k+1,\chi)}{L(k+1,\chi)}
 +\overline{\left(
       \frac{L'(-k,\chi)}{L(-k,\chi)}\right)}
 =\gamma+\log\frac{2\pi}{q}-H_k.
\end{equation}
\end{theorem}

\begin{proof}
Write $\chi(-1)=(-1)^\delta$, $\delta\in\{0,1\}$, and
\[
 \Lambda(s,\chi)=
 \left(\frac q\pi\right)^{(s+\delta)/2}
 \Gamma\left(\frac{s+\delta}{2}\right)L(s,\chi).
\]
The primitive functional equation \cite{DLMF25} implies, at real regular points,
\begin{equation}\label{bridge:completed}
 \frac{\Lambda'}{\Lambda}(s,\chi)
 =-\overline{\left(
       \frac{\Lambda'}{\Lambda}(1-s,\chi)\right)}.
\end{equation}
The positive value $L(k+1,\chi)$ is nonzero by its absolutely convergent
Euler product.  The functional equation therefore makes
$\Lambda(-k,\chi)$ finite and nonzero.  This also holds for the
conductor-$1$ character: the completed Riemann zeta function is
meromorphic, but neither $-k$ nor $k+1$ is one of its poles.

Suppose first that $\delta\equiv k\pmod2$, and put
$m=(k-\delta)/2\geq0$.  The gamma factor has a simple pole at $-k$,
so $L(s,\chi)$ must have a simple zero there.  With
$a=L'(-k,\chi)\ne0$ and $b=L''(-k,\chi)$, write $s=-k+u$ and expand:
\[
 L(-k+u,\chi)=a u+\tfrac12 b u^2+O(u^3),
\]
\[
 \Gamma(-m+u/2)=\frac{(-1)^m}{m!}
       \left(\frac2u+H_m-\gamma+O(u)\right).
\]
Taking the logarithmic derivative after cancellation gives
\begin{equation}\label{bridge:regularized-log}
 \frac{\Lambda'}{\Lambda}(-k,\chi)
 =\frac12\log\frac q\pi
  +\frac12(H_m-\gamma)+\frac{b}{2a}.
\end{equation}
At the positive point, ordinary differentiation gives
\[
 \frac{\Lambda'}{\Lambda}(k+1,\chi)
 =\frac12\log\frac q\pi
  +\frac12\psi\left(\frac{k+1+\delta}{2}\right)
  +\frac{L'(k+1,\chi)}{L(k+1,\chi)}.
\]
Insert these two formulas in \eqref{bridge:completed}.  The elementary
half-integer digamma identity \cite{DLMF5}
\[
 \psi\left(\frac{k+1+\delta}{2}\right)+H_m-\gamma
 =2H_k-2\gamma-2\log2
\]
follows from $\psi(r+\tfrac12)=-\gamma-2\log2+2H_{2r}-H_r$,
checking $\delta=0$ and $\delta=1$.  It proves
\eqref{bridge:formula}, including the factor $1/2$.

In the opposite parity the gamma factor is finite and nonzero at
$-k$, hence $L(-k,\chi)\ne0$.  Applying
\eqref{bridge:completed} directly gives the sum of the two logarithmic
$L$-derivatives as
\[
 -\log\frac q\pi
 -\frac12\psi\left(\frac{k+1+\delta}{2}\right)
 -\frac12\psi\left(\frac{\delta-k}{2}\right).
\]
The second digamma argument is a half-integer.  Digamma reflection
there has zero cotangent term, and duplication reduces the sum to
$2H_k-2\gamma-2\log2$, proving
\eqref{bridge:nontrivial}.
\end{proof}

For example, the trivial-zero cases $\zeta(-2)=0$ and $\beta(-1)=0$
give respectively
\[
 \frac{\zeta'(3)}{\zeta(3)}
 +\frac{\zeta''(-2)}{2\zeta'(-2)}
 =\gamma+\log(2\pi)-\frac32,
\]
\[
 \frac{\beta'(2)}{\beta(2)}
 +\frac{\beta''(-1)}{2\beta'(-1)}
 =\gamma+\log\frac\pi2-1.
\]
Thus the positive-side first derivative really corresponds to a
negative-side second derivative at the trivial-zero parity.  It cannot
in general be recovered from the negative first-jet value alone.
The accompanying numerical check includes these two identities and
complex primitive quartic characters modulo $5$, for which the
conjugation in \eqref{bridge:formula} is essential.

% ----- sections/gaussian_reduction.tex -----
\section{A proved reduction at every even weight in the Gaussian ladder}
\label{gau:section}

The draft \emph{Gaussian multiple polylogarithm depth} studies
\begin{equation}\label{gau:sum-definition}
 S_p=\sum_{n\geq 0}\frac{(-1)^nH_n}{(2n+1)^p},
 \qquad H_n=\sum_{j=1}^n\frac1j,\quad H_0=0.
\end{equation}
Its formulas for $S_5$ and $S_7$ are correct, but the argument that a new
Dirichlet beta value ``absorbs'' the depth-two content is not a proof of a
reduction at every even weight. Here we supply a complete analytic proof and
an exact generating function.

\paragraph{Prior art and the precise contribution.}
The all-weight formula below already appears in a 2022 answer by
Ali Olaikhan, including exactly the $S_5$ and $S_7$ specializations displayed
in the repository draft \cite{gau-olaikhan-2022}.
It therefore requires attribution rather than a novelty claim. The purpose
of this section is to give a self-contained proof with convergence and
parameter domains specified, connect the formula to the draft's depth
question, and extract a convergent resummation. None of these arguments
asserts irrationality, algebraic independence, or minimal depth.

\subsection{Integral representation and a reflection identity}

For every integer $p\geq1$, the sum in \eqref{gau:sum-definition} converges.
For $p\geq2$ it converges absolutely. For $p=1$, the positive sequence
$H_n/(2n+1)$ decreases to zero: indeed,
\[
 \frac{H_{n+1}}{2n+3}\leq\frac{H_n}{2n+1}
 \quad\Longleftrightarrow\quad
 \frac{2n+1}{n+1}\leq2H_n,
\]
which holds for $n\geq1$. The alternating-series test applies.

\begin{lemma}\label{gau:moment}
For every integer $p\geq1$,
\begin{equation}\label{gau:moment-formula}
 S_p=-\frac1{(p-1)!}\int_0^1
       \frac{(-\log x)^{p-1}\log(1+x^2)}{1+x^2}\,dx.
\end{equation}
\end{lemma}
\begin{proof}
Summing the geometric series after the inner harmonic-number sum gives
\[
 \sum_{n\geq0}H_n y^n=-\frac{\log(1-y)}{1-y},\qquad |y|<1.
\]
First use $y=-r x^2$ with $0<r<1$, multiply by
$(-\log x)^{p-1}/(p-1)!$, and integrate term by term. The elementary
Mellin integral of $x^{2n}$ gives $(2n+1)^{-p}$.
Now let $r\uparrow1$. Abel's theorem applies to the convergent alternating
series. On the integral side,
$\log(1+rx^2)/(1+rx^2)$ is uniformly bounded on $0\leq x\leq1$,
and $(-\log x)^{p-1}$ is integrable, so dominated convergence applies.
This also covers the conditionally convergent case $p=1$.
\end{proof}

Introduce
\begin{align}
 J(s)&=\int_0^1\frac{x^{s-1}\log(1+x^2)}{1+x^2}\,dx,
 & C(s)&=\int_0^1\frac{x^{s-1}}{1+x^2}\,dx,\label{gau:mellin-defs}\\
 K(s)&=\int_0^\infty\frac{x^{s-1}\log(1+x^2)}{1+x^2}\,dx.
 \nonumber
\end{align}
The functions $J$ and $C$ are holomorphic for $\Re s>-2$ and $\Re s>0$,
respectively; $K$ is holomorphic for $-2<\Re s<2$.
In the common strip $0<\Re s<2$, the substitution $x=1/t$ in the
part of $K$ above $1$ gives
\begin{equation}\label{gau:reflection}
 K(s)=J(s)+J(2-s)-2C'(2-s).
\end{equation}
All signs in this identity matter: the logarithm transforms as
$\log(1+t^{-2})=\log(1+t^2)-2\log t$.

The beta integral gives another evaluation. For
$0<\Re s<2\Re a$,
\[
 \int_0^\infty\frac{x^{s-1}}{(1+x^2)^a}\,dx
 =\frac{\Gamma(s/2)\Gamma(a-s/2)}{2\Gamma(a)}.
\]
Taking minus the derivative in $a$ at $a=1$ yields
\begin{equation}\label{gau:complete-mellin}
 K(s)=\frac{\pi}{2\sin(\pi s/2)}
          \bigl[\psi(1)-\psi(1-s/2)\bigr],
 \qquad 0<\Re s<2.
\end{equation}
Differentiation under the integral is justified locally uniformly in that
strip by integrable power bounds with a logarithmic factor.

\subsection{The finite reduction}

Write
\[
 \beta(r)=\sum_{n\geq0}\frac{(-1)^n}{(2n+1)^r},
 \qquad
 \sec t=\sum_{j\geq0}\frac{|E_{2j}|}{(2j)!}t^{2j},
\]
so that $|E_0|,|E_2|,|E_4|,|E_6|,|E_8|=1,1,5,61,1385$.
Define $b_0=\log2$ and
$b_k=(1-2^{-2k-1})\zeta(2k+1)$ for $k\geq1$.

\begin{theorem}[All even-weight reductions]\label{gau:all-weights}
For every integer $m\geq1$,
\begin{equation}\label{gau:finite-reduction}
 \boxed{\displaystyle
 S_{2m-1}=(2m-1)\beta(2m)
 -\frac\pi2\sum_{k=0}^{m-1}
 \frac{|E_{2m-2-2k}|}{(2m-2-2k)!}
       \left(\frac\pi2\right)^{2m-2-2k}b_k.}
\end{equation}
Consequently $S_{2m-1}$ is a rational linear combination of
$\beta(2m)$, $\pi^{2m-1}\log2$, and
$\pi^{2m-2k-1}\zeta(2k+1)$, $1\leq k\leq m-1$.
\end{theorem}
\begin{proof}
For $|z|<1$, take the even part of \eqref{gau:reflection} about $s=1$.
If $K_{\mathrm{ev}}(z)=(K(1+z)+K(1-z))/2$, then
\begin{equation}\label{gau:even-reflection}
 J(1+z)+J(1-z)
 =K_{\mathrm{ev}}(z)+C'(1+z)+C'(1-z).
\end{equation}
The moment representation implies
\[
 J(1+z)+J(1-z)=-2\sum_{m\geq1}S_{2m-1}z^{2m-2}.
\]
Differentiating $C$ under its integral, then using the geometric series,
gives
\[
 C'(1+z)+C'(1-z)
 =-2\sum_{m\geq1}(2m-1)\beta(2m)z^{2m-2}.
\]
These expansions converge locally uniformly for $|z|<1$.

Next, \eqref{gau:complete-mellin} yields
\[
 K_{\mathrm{ev}}(z)=\frac\pi{2\cos(\pi z/2)}
 \left[\psi(1)-\frac{\psi((1-z)/2)+\psi((1+z)/2)}2\right].
\]
The digamma series, obtained by termwise differentiation of
$\psi(w)=-\gamma+\sum_{n\geq0}((n+1)^{-1}-(n+w)^{-1})$, gives
\[
 \psi(1)-\psi(\tfrac12)=2\log2,
 \qquad
 \psi^{(2k)}(\tfrac12)=-(2k)!(2^{2k+1}-1)\zeta(2k+1)
 \quad(k\geq1).
\]
It follows that
\[
 K_{\mathrm{ev}}(z)
 =\pi\sec(\pi z/2)\sum_{k\geq0}b_k z^{2k}.
\]
Substitute these three series into \eqref{gau:even-reflection} and compare
the coefficient of $z^{2m-2}$. The result is
\eqref{gau:finite-reduction}.
\end{proof}

For orientation, the first two instances are
\[
 S_1=\beta(2)-\frac\pi2\log2,
 \qquad
 S_3=3\beta(4)-\frac{\pi^3}{16}\log2
                     -\frac{7\pi}{16}\zeta(3).
\]
The cases $m=3,4$ give exactly the draft's equations
\texttt{eq:S5} and \texttt{eq:S7}. The next instance is
\begin{align}\label{gau:s9}
 S_9={}&9\beta(10)-\frac{277\pi^9}{4128768}\log2
       -\frac{427\pi^7}{737280}\zeta(3)
       -\frac{155\pi^5}{24576}\zeta(5)\nonumber\\
     &\hspace{14mm}-\frac{127\pi^3}{2048}\zeta(7)
       -\frac{511\pi}{1024}\zeta(9).
\end{align}
Every term has total weight $10$.

\paragraph{What parity does and does not prove.}
The mechanism is the reflection $s\mapsto2-s$ of the Mellin integral.
Its even Taylor part isolates the odd powers of the denominator. This
proves the upper bound on the complexity of $S_{2m-1}$ without assuming
that any even beta value is a new transcendental.
It gives no lower bound for the complexity of $S_{2m}$.
In particular, neither the collapse of $\beta(2m+1)$ to a rational multiple
of $\pi^{2m+1}$ nor a failed integer-relation search proves that
$S_{2m}$ requires a double polylogarithm.

\subsection{An exact generating function and a resummation}

\begin{proposition}\label{gau:generating}
The generating function
$F(z)=\sum_{m\geq1}S_{2m-1}z^{2m-2}$ has Taylor radius exactly $3$.
It is represented, for $|\Re z|<3$, by
\begin{equation}\label{gau:generating-integral}
 F(z)=-\int_0^1\frac{\cosh(z\log x)\log(1+x^2)}{1+x^2}\,dx.
\end{equation}
For $|z|<1$ it also has the explicit form
\begin{equation}\label{gau:generating-digamma}
\begin{split}
 F(z)={}&-\frac{C'(1+z)+C'(1-z)}2\\
 &-\frac{\pi}{4\cos(\pi z/2)}
 \left[\psi(1)-\frac{\psi((1-z)/2)+\psi((1+z)/2)}2\right],
\end{split}
\end{equation}
where
$C'(s)=\bigl[\psi'((s+2)/4)-\psi'(s/4)\bigr]/16$.
The apparent singularities of the whole right side at $z=\pm1$ are removable.
Moreover,
\begin{equation}\label{gau:resummation}
 \boxed{\displaystyle\sum_{m\geq1}S_{2m-1}=-\frac{\pi^2}{48}.}
\end{equation}
\end{proposition}
\begin{proof}
The alternating-series estimate gives $|S_p|\leq3^{-p}$ for $p\geq1$.
It follows that the Taylor series converges for $|z|<3$, and summing the
moments in \eqref{gau:moment-formula} gives
\eqref{gau:generating-integral}. The integrand is dominated on compact
subsets of $|\Re z|<3$ because $\log(1+x^2)/(1+x^2)=O(x^2)$ at zero.
Equation \eqref{gau:generating-digamma} follows from
\eqref{gau:even-reflection}. Its apparent singularities at $\pm1$ must
cancel by \eqref{gau:generating-integral}.

More explicitly, subtracting the constant coefficient gives the meromorphic
continuation
\begin{equation}\label{gau:partial-fractions}
 F(z)=S_1+\sum_{n\geq1}
 \frac{(-1)^nH_n z^2}{(2n+1)((2n+1)^2-z^2)}.
\end{equation}
The series on the right is absolutely and locally uniformly convergent
away from $z=\pm(2n+1)$, $n\geq1$. Its $n=1$ summand has a nonzero simple
pole at $z=3$, while all remaining summands are analytic there. Therefore
the Taylor radius is exactly $3$.
Finally, setting $z=1$ in \eqref{gau:generating-integral} simplifies the
kernel:
\[
 F(1)=-\frac12\int_0^1\frac{\log(1+x^2)}x\,dx
      =-\frac14\int_0^1\frac{\log(1+t)}t\,dt
      =-\frac{\pi^2}{48}.
\]
The last integral follows from the absolutely convergent integrated
alternating series for $\log(1+t)$.
\end{proof}

\paragraph{Reproducible checks.}
The program \texttt{code/verify\_gaussian.py} computes $p=1,3,5,7,9$ in
three ways: accelerated summation of \eqref{gau:sum-definition}, numerical
quadrature of \eqref{gau:moment-formula}, and the finite expression
\eqref{gau:finite-reduction}. At $90$ decimal digits the largest pairwise
absolute residual was less than $4\cdot10^{-91}$. It also checks five
values of the generating function, including a complex argument and
$z=3/2$, and the removable-singularity value $F(1)$.
The values and precision settings are recorded in
\texttt{data/gaussian\_checks.json}. These checks detect implementation
and transcription errors; the proofs above establish the identities.

% ----- sections/audit_introduction.tex -----
\section{Corrections and qualifications for the existing corpus}
\label{audit:section}

The following findings refer to the pinned source snapshot specified
in Section~\ref{intro:section}. Article filenames are relative to
\path{Analysis/Polylogarithms/docs/articles}; report filenames are
relative to \path{Analysis/Polylogarithms/docs/reports}.
Source labels are given where available, so that later line-number
changes do not obscure the intended location.
The separate \path{CORRECTIONS.md} register organizes the integration
actions and records which issues the intake README already identified.

Three different responses are needed. A false formula or false
classification requires a corrected statement and a witness. A true
statement with an incomplete argument requires a proof, as supplied
for the rank laws in Section~\ref{rank:section}. A claim of arithmetic
independence requires an appropriately stated conjecture or a new
independence theorem; numerical searches do not settle it.
The corrections below preserve the proved evaluations while specifying
which conclusions are justified by their proofs.

% ----- sections/corrections_polylogs.tex -----
\subsection{Absolute convergence at the boundary}
\label{audit:poly:convergence}

\paragraph{Source.}
\texttt{multiple-polylogarithms-introduction.tex}, the theorem titled
``Domain of absolute convergence'' (source lines 231--235).
The condition printed there, $\Re s_j\geq1$, $|z_j|\leq1$, and
$(s_1,z_1)\neq(1,1)$, does not imply absolute convergence.
For example, $\Li_1(i)$ satisfies every stated hypothesis, while the
absolute-value series is $\sum_{n\geq1}1/n$.

\paragraph{Replacement.}
For $\Re s_j\geq1$ and $|z_j|\leq1$, a sufficient condition for
absolute convergence is
\[
 |z_1|<1\qquad\text{or}\qquad\Re s_1>1.
\]
Indeed, the absolute value of the inner nested sum is at most
$H_{n_1-1}^{k-1}/(k-1)!$, so the full series is dominated by
\[
 \frac1{(k-1)!}\sum_{n\geq1}
    \frac{|z_1|^n(1+\log n)^{k-1}}{n^{\Re s_1}}.
\]
Conditional convergence at boundary roots of unity is a separate issue.
The usual positive-integer MZV admissibility condition $s_1\geq2$ remains
correct. Also, the definition of $S_p$ in
\texttt{gaussian-multiple-polylog-depth.tex} should start at $p\geq1$:
its printed allowance $p=0$ gives terms $(-1)^nH_n$ that do not tend to
zero. An Abel-regularized value would need its own explicit definition.

\subsection{The order of the iterated-integral letters}
\label{audit:poly:integral-order}

\paragraph{Source.}
\texttt{gaussian-eisenstein-double-polylogs.tex},
\texttt{eq:G-def} and \texttt{eq:goncharov-dbl}; the same mismatch appears
in \texttt{multiple-polylogarithms-introduction.tex}, lines 276--282.
The displayed definition attaches the letter $a_i$ to $t_i$ on the simplex
$0<t_1<\cdots<t_n<1$, but the subsequent polylogarithm dictionary uses the
opposite convention. In particular, the displayed integral for
$G(0,z^{-1};1)$ contains $\int_0^{t_2}dt_1/t_1$, which diverges, although
$-\Li_2(z)$ is finite, for example at $z=1/2$.

\paragraph{Replacement.}
Keep the stated dictionary and define the leftmost letter to be outermost:
\begin{align*}
 G(a_1,\ldots,a_n;z)
   &=\int_0^z\frac{G(a_2,\ldots,a_n;t)}{t-a_1}\,dt,\\
 G(a_1,\ldots,a_n;1)
   &=\int_{0<t_n<\cdots<t_1<1}
            \prod_{j=1}^n\frac{dt_j}{t_j-a_j}.
\end{align*}
For convergent words this gives
\[
 \Li_n(z)=-G(0^{n-1},z^{-1};1),\qquad
 \Li_{a,b}(u,v)=G(0^{a-1},u^{-1},0^{b-1},(uv)^{-1};1).
\]
Alternatively, one may retain the increasing-index simplex and reverse
every letter word in the dictionary. The two choices must not be mixed.
Endpoint regularization, when needed, should be stated separately.

\subsection{An extra term in the modified polylogarithm}
\label{audit:poly:modified}

\paragraph{Source.}
\path{ladders-as-bloch-elements__ce20ea60a460.tex},
\texttt{eq:Pm}, and the reported complex-embedding value of $P_4$.
The upper endpoint of the sum is printed as $m$. For $m\geq2$ the correct
modified polylogarithm is
\begin{equation}\label{audit:poly:correct-P}
 P_m(z)=\operatorname{Re}_m\!\left(
       \sum_{j=0}^{m-1}\frac{2^jB_j}{j!}
          \log^j|z|\,\Li_{m-j}(z)\right),
 \quad
 \operatorname{Re}_m=
 \begin{cases}\Re,&m\ \text{odd},\\\Im,&m\ \text{even}.
 \end{cases}
\end{equation}
Here $B_1=-1/2$. This agrees with the standard $P_m$ definition in the
PARI/GP documentation, \texttt{polylog} with flag $3$
\cite{audit-pari-polylog}; its additional real multiple of
$B_m\log^m|z|$ vanishes after $\operatorname{Re}_m$ for $m\geq2$.

The printed expression adds, at even $m$, the generally nonzero term
\begin{equation}\label{audit:poly:P-error}
 \frac{2^mB_m}{m!}\log^m|z|\,
             \Im\!\left(\frac{z}{1-z}\right).
\end{equation}
For example, at $m=2$, $z=i/2$, the excess is exactly
$2\log^2(2)/15$. Verification only at odd weights or at $z=i$ cannot
detect this error, because respectively $B_m=0$ or $\log|z|=0$.

The error also affects a numerical conclusion. For the lower nonreal
root of $z^3+z-1=0$,
\[
z_-=-0.34116390191400966\ldots-1.16154139999725194\ldots i,
\]
recomputation gives
\begin{align*}
 P_4(z_-)&=-0.916010467826680838722383200166\ldots,\\
 P_4^{\mathrm{printed}}(z_-)&=-0.915999527027516872882916723105\ldots.
\end{align*}
The latter agrees with the report's $-0.91599\ldots$.
Thus the definition and the complex-embedding computation both require
correction; the odd-weight formulas are not affected by this particular
endpoint error.

\subsection{The asserted binomial dimensions are not dimensions of depth}
\label{audit:poly:depth-dimensions}

\paragraph{Source.}
\texttt{gaussian-eisenstein-double-polylogs.tex},
\texttt{thm:deligne}, \texttt{eq:binomial-dim}, and
\texttt{cor:depth3}. The formula
$\dim\operatorname{gr}^D_p(\text{weight }n)=\binom np$
is false for the ordinary cyclotomic depth filtration used in the cited
literature.

A direct contradiction already occurs at $n=6$, $p=1$. All level-four
weight-six single polylogarithms lie in the span of $\zeta(6)$ and
$i\beta(6)$, since
\[
 \Li_6(-1)=-\frac{31}{32}\zeta(6),\qquad
 \Li_6(\pm i)=-\frac{31}{2048}\zeta(6)\pm i\beta(6).
\]
The depth-one space therefore has dimension at most $2$, whereas the
printed formula gives $6$. Likewise, the depth-zero component of positive
weight is zero, not $\binom n0=1$.
These contradictions use explicit identities and do not depend on
unproved numerical-period independence.

Glanois defines depth by the number of nonzero iterated-integral letters
and explicitly warns that it is a filtration rather than a grading on
$\mathcal H$ \cite[\S3.1]{audit-glanois-2016}.
A Hilbert series in weight alone does not determine these filtered
pieces. Moreover, even a correct positive dimension for an ambient
motivic depth-three space would not imply that three particular
$\Li_{a,b,c}(i,1,1)$ have nonzero independent images in it.
An explicit computation of their classes is required.
The period conjecture does not repair either of these logical gaps.

\paragraph{Replacement.}
Withdraw \texttt{eq:binomial-dim} and the stated proof of
\texttt{cor:depth3}. Retain the verified shuffle relations as reductions
modulo their displayed product terms. The minimal depth of the three
specified numerical survivors remains undetermined by the arguments in
that draft. A motivic claim should name the precise filtration and supply
the relevant coaction or basis computation before making a corresponding
conditional numerical claim.

\subsection{Rank seven does not mean seven individual triples collapse}
\label{audit:poly:shuffle-rank}

\paragraph{Source.}
\texttt{gaussian-eisenstein-double-polylogs.tex},
\texttt{prop:7of10} and its preceding rank calculation.
There are ten compositions $(a,b,c)$ of weight $6$ with positive entries.
The draft obtains rank $7$ from shuffle products and concludes that seven
of the ten triples are themselves products. The correct conclusion is
that seven independent \emph{linear combinations} are products.

We recomputed the $13$ formal shuffle-product rows exactly: the ten
products of a single word and a double word, together with the three
unordered triple products of single words. Their rational rank is indeed
$7$. However, none of the ten individual coordinate vectors belongs to
the row space of these products.
This can be checked using the following three integer linear functionals,
in lexicographic order of the compositions
$(1,1,4),(1,2,3),\ldots,(4,1,1)$:
\[
 \begin{pmatrix}
 -1&2&-2&1&0&0&-1&1&0&0\\
 -3&6&-6&3&-1&4&-4&0&1&0\\
 -6&12&-12&6&-3&12&-9&0&0&1
 \end{pmatrix}.
\]
Each row annihilates every shuffle-product row, the three rows are
independent, and each of the ten columns is nonzero.
Thus every individual word is detected by at least one annihilating
functional. The full matrix and exact checks are generated by
\texttt{code/verify\_shuffle\_audit.py} and recorded in
\texttt{data/gaussian\_shuffle\_audit.json}.

\paragraph{Replacement.}
The formal weight-six shuffle quotient modulo products is
three-dimensional, and seven of its ten coordinate variables can be
eliminated \emph{in terms of three survivors and products}.
This formal calculation neither proves nor disproves additional
relations after evaluation at $i$.

The same distinction applies to \texttt{thm:gauss-5} and
\texttt{thm:eis-1}: a corank computed from a specified relation matrix is
an upper bound on the dimension of the corresponding numerical span,
unless completeness of those relations is independently established.

\subsection{Two further logical boundary corrections}
\label{audit:poly:other}

In \texttt{gaussian-multiple-polylog-depth.tex}, the discussion of the
weight-four symmetric triple sum compares it to
$\zeta(2,1,1)+\zeta(1,2,1)+\zeta(1,1,2)$ as though this were a convergent
uncolored MZV combination. The last two defining sums diverge. In each,
the inner sum is bounded below by a positive constant for all sufficiently
large outer indices, and the outer exponent is $1$; comparison with the
harmonic series proves divergence. The comparison should be removed,
or replaced by an explicitly specified regularized identity.

The same draft infers at least $50$ independent new directions from
$50$ individual failures of membership in a fixed span. Even exact
non-membership of each value would not justify that count: all $50$
values could represent one common nonzero class in the quotient.
Numerical failures of membership provide still less. The appropriate
replacement is a list of tested candidate values with no claimed lower
bound on the dimension of their joint span.

% ----- sections/corrections_gamma.tex -----
\subsection{Log-gamma integrals: the missing even-character contribution}
\label{corr:loggamma}

The abstract and Section~5 of
\path{polylog-polygamma-bridge.tex} suggest that
$\psi^{(-2)}(a)=\int_0^a\log\Gamma(t)\,dt$ at every rational $a$ reduces
to weight-two Clausen values, $\zeta'(-1)$, and logarithms.  The three
displayed examples, at denominators $3,4,6$, are correct.  Their
extrapolation to all denominators omits an even-character contribution.
The same extrapolation occurs in Section~1 of
\path{loggamma-integrals-clausen-bridge__c5013328db83.md}.
Here is an explicit correction at denominator $5$.

\begin{proposition}[An exact denominator-five correction]
\label{corr:loggamma-five}
Let $\chi_5$ be the quadratic character modulo $5$, with values
$(1,-1,-1,1)$ on residues $1,2,3,4$.  Then
\begin{align}
 \int_0^{1/5}\log\Gamma(t)\,dt
 ={}&\frac{2}{25}+\frac{\log(2\pi)}{10}
       +\frac{\operatorname{Cl}_2(2\pi/5)}{4\pi}
       -\frac65\zeta'(-1) \notag\\
    &+\frac1{20}L'(-1,\chi_5)-\frac{29}{1200}\log5.
 \label{corr:loggamma-five-formula}
\end{align}
More generally, for $a>0$,
\begin{equation}\label{corr:loggamma-general}
 \int_0^a\log\Gamma(t)\,dt
 =\zeta'(-1,a)-\zeta'(-1)
       +\frac{a-a^2}{2}+\frac a2\log(2\pi).
\end{equation}
\end{proposition}
\begin{proof}
Differentiate $\partial_a\zeta(s,a)=-s\zeta(s+1,a)$ with respect to
$s$ at $s=-1$.  Lerch's identity and $\zeta(0,a)=\tfrac12-a$ give
\[
 \partial_a\zeta'(-1,a)
 =\zeta'(0,a)-\zeta(0,a)
 =\log\Gamma(a)-\tfrac12\log(2\pi)-\tfrac12+a.
\]
The shift identity for Hurwitz zeta shows that
$\zeta'(-1,a)\to\zeta'(-1)$ as $a\downarrow0$, since the additional
term is $-a\log a$.  Integration proves
\eqref{corr:loggamma-general}.  These are classical Hurwitz-zeta
identities \cite{DLMF25}; the integrated theory is also developed by
Espinosa and Moll \cite{EspinosaMoll2002b}.

Put $Z_r=\zeta'(-1,r/5)$ and $Z=\zeta'(-1)$.  The Fourier reflection
formula at this weight is
\begin{equation}\label{corr:zeta-minus-one-reflection}
 \zeta'(-1,a)-\zeta'(-1,1-a)
       =\frac{\operatorname{Cl}_2(2\pi a)}{2\pi},
       \qquad 0<a<1.
\end{equation}
For completeness, subtract the two Hurwitz Fourier expansions
\[
 \zeta(1-s,a)-\zeta(1-s,1-a)
 =\frac{4\Gamma(s)}{(2\pi)^s}\sin\frac{\pi s}{2}
       \sum_{m\geq1}\frac{\sin(2\pi ma)}{m^s}
\]
and differentiate at $s=2$.  Absolute local convergence of the
series and its logarithmic derivative justifies the differentiation;
the zero of the sine factor leaves precisely
\eqref{corr:zeta-minus-one-reflection}.

Differentiating distribution and the character decomposition at $s=-1$
gives
\begin{align*}
 Z_1+Z_2+Z_3+Z_4&=-\frac45 Z-\frac{\log5}{60},\\
 Z_1-Z_2-Z_3+Z_4
   &=\frac15\bigl(L'(-1,\chi_5)+\log5\,L(-1,\chi_5)\bigr).
\end{align*}
Here $L(-1,\chi_5)=-2/5$, obtained directly from
$\zeta(-1,a)=-B_2(a)/2$.  Combining these equations with
$Z_1-Z_4=\operatorname{Cl}_2(2\pi/5)/(2\pi)$ yields
\[
 Z_1=-\frac15 Z+\frac1{20}L'(-1,\chi_5)
       -\frac{29}{1200}\log5
       +\frac{\operatorname{Cl}_2(2\pi/5)}{4\pi}.
\]
Insert this into \eqref{corr:loggamma-general} at $a=1/5$.
\end{proof}

The correction concerns the scope of the available reduction laws.
It does not prove that $L'(-1,\chi_5)$ is transcendental or independent
of the other constants in \eqref{corr:loggamma-five-formula}.  A precise
replacement for the earlier assertion is: the odd reflection component
is Clausen-explicit, while the even component requires the appropriate
Dirichlet-$L$ derivatives unless further relations have been proved.
The later \path{stieltjes-antiderivative-ladder.tex} already gives the
general Hurwitz-jet description; the earlier bridge discussion should
be brought into agreement with it.

\subsection{Two Herglotz classification statements have exact counterexamples}
\label{corr:herglotz}

In this subsection $F$ denotes the Herglotz--Zagier function, and
$J$ its integral companion:
\[
 F(x)=\sum_{m\geq1}\frac{\psi(mx)-\log(mx)}{m},\qquad
 J(x)=\int_0^1\frac{\log(1+t^x)}{1+t}\,dt,\qquad x>0.
\]
Their classical relation is
\begin{equation}\label{corr:J-F}
 J(x)=F(2x)-2F(x)+F(x/2)+\frac{\pi^2}{12x}.
\end{equation}
See Radchenko and Zagier \cite{RadchenkoZagier2023} for these functions,
the functional equations, and the finite dilogarithm representation.
They are unrelated to the function $F_{n,k}$ in
Section~\ref{zero:section}.

The proposition labelled \texttt{prop:collapse} in
\path{herglotz-arithmetic-study__467d9018ec87.tex} asserts a criterion
involving $\operatorname{rad}(q)\mid30$ for reducing $F(p/q)-F(1)$
to rational-argument dilogarithms, logarithms, and $\pi^2$.
Its necessity fails already at $p/q=1/7$.  Indeed, the report's own
proved formula \texttt{thm:F1q} is
\begin{equation}\label{corr:F-unit-fraction}
 F(1/q)-F(1)
 =-\frac{(q-1)(3q-2)}{24q}\pi^2-\frac12 S(q),
 \qquad
 S(q)=\sum_{r=1}^{q-1}\log^2\left(2\sin\frac{\pi r}{q}\right).
\end{equation}
In particular,
\begin{equation}\label{corr:F-seventh}
 F(1/7)-F(1)
   =-\frac{19}{28}\pi^2
      -\sum_{r=1}^{3}\log^2\left(2\sin\frac{\pi r}{7}\right).
\end{equation}
This contains no dilogarithm at all, although $7\nmid30$.

To check the mechanism independently, specialize the finite
root-of-unity dilogarithm formula to numerator $1$:
\[
 F(1/q)-F(1)=
 \sum_{\substack{\beta^q=1\\\beta\ne1}}
   \left[\operatorname{Li}_2\left(\frac{\beta}{\beta-1}\right)
          -\frac{\pi^2}{6}\right].
\]
If $\beta=e^{2\pi i r/q}$, the transformed argument has real part
$1/2$, modulus $(2\sin(\pi r/q))^{-1}$, and argument
$\pi(r/q-1/2)$.  Euler reflection, paired with conjugation, evaluates
its real dilogarithm component.  Summing the squared arguments and
treating $\beta=-1$ separately when $q$ is even gives
\eqref{corr:F-unit-fraction}.

Nor does restricting to numerators greater than $1$ repair the criterion.
The same laws give
\begin{align}
 F(6/7)-F(1)={}&-\frac8{63}\pi^2+\operatorname{Li}_2(6/7)
       +\frac12\log^2(7/6)\notag\\
       &-\frac12S(7)+\frac12S(6).
 \label{corr:F-six-sevenths}
\end{align}
For a direct derivation, apply the three-term functional equation to
$x=1/6$ and the inversion formula to $x=7/6$:
\begin{align*}
 F(x)-F(x+1)-F\left(\frac{x}{x+1}\right)
   &=-F(1)+\operatorname{Li}_2\left(\frac1{x+1}\right),\\
 F(x)+F(1/x)-2F(1)
   &=\frac12\log^2x-\frac{\pi^2(x-1)^2}{6x}.
\end{align*}
Eliminate $F(7/6)$ and use \eqref{corr:F-unit-fraction} for $q=6,7$.
Thus the appropriate correction is to remove the asserted
classification and retain the proved families and individually
established reductions.  These counterexamples do not determine a
replacement classification for all numerators and denominators.

The proposition labelled \texttt{prop:J} in the same report asserts
that $J(2/q)$ has the specified elementary form only for $q=3,5$.
The allowed elementary inventory in its examples includes $\pi^2$.
With that convention, $q=4$ is another exact counterexample:
\begin{equation}\label{corr:J-half}
 \boxed{\displaystyle J(1/2)=\frac{\pi^2}{48}
                                  +\frac14\log^22.}
\end{equation}
Indeed, \eqref{corr:F-unit-fraction} gives
\[
 F(1/2)-F(1)=-\frac{\pi^2}{12}-\frac12\log^22,
 \qquad
 F(1/4)-F(1)=-\frac{5\pi^2}{16}-\frac34\log^22.
\]
Their substitution in \eqref{corr:J-F} proves \eqref{corr:J-half}.
It also disproves the following sentence in that proposition, which
claims that $J(1/q)$ always retains dilogarithms.  Adding an odd-$q$
restriction would avoid this particular example, but would not prove
the remaining necessity claim.

\subsection{The even-denominator term in the Herglotz derivative formula}
\label{corr:herglotz-derivative}

The same report's \texttt{thm:deriv} defines
\[
 W(q)=\sum_{r=1}^{q-1}
       \cot\frac{2\pi r}{q}\,
       \operatorname{Cl}_2\left(\frac{2\pi r}{q}\right).
\]
For even $q$ the displayed summand at $r=q/2$ is undefined as written:
$\cot\pi$ is singular and $\operatorname{Cl}_2(\pi)=0$.
Define instead
\begin{equation}\label{corr:W-definition}
 W_*(q)=\sum_{\substack{1\leq r<q\\2r\ne q}}
       \cot\frac{2\pi r}{q}\,
       \operatorname{Cl}_2\left(\frac{2\pi r}{q}\right).
\end{equation}
The corrected formula, retaining the report's separate correction, is
\begin{equation}\label{corr:F-derivative}
 \frac2qF'(2/q)-(1+\log2)
    =\frac{q^2-4}{24q}\pi^2+W_*(q)
                                  -\mathbf1_{2\mid q}\log2.
\end{equation}
Equivalently, extend the product continuously at angle $\pi$ and
include it in the sum, while removing the separate $-\log2$ term.
These conventions must not be combined.

The continuous value follows from
$\operatorname{Cl}_2'(\theta)=-\log(2\sin(\theta/2))$ near $\pi$:
\[
 \operatorname{Cl}_2(\pi+u)=-u\log2+O(u^3),\qquad
 \cot(\pi+u)=u^{-1}+O(u),
\]
so the product tends to $-\log2$.  This explains exactly the finite
term required in \eqref{corr:F-derivative}.
At $q=4$ both remaining nonzero-angle cotangents vanish, so
$W_*(4)=0$ and
\begin{equation}\label{corr:F-derivative-half}
 F'(1/2)=2+\frac{\pi^2}{4}.
\end{equation}
Consequently the accompanying statement that this derivative family
never reduces to elementary form is also too broad.

\subsection{Cubic subfields, class fields, and the regulator coefficient}
\label{corr:stark-fields}

The report
\path{herglotz-stark-regulators__6096eccfa2a0.tex}, in its discussion of
$\mathbb Q(\sqrt{257})$ and in the subsequent table, uses the same
symbol $H_0$ both for a cubic field over $\mathbb Q$ and for a cyclic
cubic class-field extension of a quadratic field.  These have different
absolute degrees.  A degree-three field over $\mathbb Q$ cannot
contain a quadratic field, by the tower law.

In the intended $S_3$ situation the notation should distinguish
\[
 [H:K]=3,\quad[K:\mathbb Q]=2,\quad[H:\mathbb Q]=6,
 \qquad [H_0:\mathbb Q]=3.
\]
Here $H/K$ is the cyclic cubic Hilbert or ring class field extension,
and $H_0$ is a non-Galois cubic subfield of its normal closure $H$.
The polynomial displayed for $H_0$ specifies the latter field, not
$H/K$.  This is the standard class-field/cubic-subfield correspondence
\cite{MilneCFT}; it is compatible with the framework of
\cite{RadchenkoZagier2023}.

There is a related omitted factor in the general explanation of the
leading $L$-coefficient.  Let $\chi$ be either nontrivial character of
$\operatorname{Gal}(H/K)\simeq C_3$.  The two conjugate induced
characters are the standard two-dimensional representation of $S_3$.
Artin formalism therefore gives
\begin{equation}\label{corr:artin-cubic}
 L_K(s,\chi)=\frac{\zeta_{H_0}(s)}{\zeta(s)}.
\end{equation}
One can see the representation identity directly: the permutation
representation on the three embeddings of $H_0$ is the direct sum
of the trivial representation and that standard representation.
If $H_0$ is totally real, the analytic class number formula at zero
then gives \cite{MilneCFT}
\[
 \zeta_{H_0}(s)=-\frac{h(H_0)\operatorname{Reg}(H_0)}{2}s^2+O(s^3),
 \qquad \zeta(0)=-\frac12,
\]
and hence
\begin{equation}\label{corr:leading-cubic-L}
 L_K^*(0,\chi):=\lim_{s\to0}s^{-2}L_K(s,\chi)
            =h(H_0)\operatorname{Reg}(H_0).
\end{equation}
Thus the assertion $L_K^*(0,\chi)=\operatorname{Reg}(H_0)$ needs
$h(H_0)=1$.  The printed general explanation omits that hypothesis.
For the five cubics actually displayed, however, it can be supplied
by a short exact argument.  This repairs the stated leading-coefficient
formula for these fields while distinguishing it from the general
identity \eqref{corr:leading-cubic-L}.  Neither this argument nor the
identity proves the report's separate fitted Herglotz evaluations.
The leading-coefficient identity follows from Artin formalism and the
analytic class number formula; it does not need a new appeal to a
conjectural Stark statement.

\begin{proposition}[Class-number certificates for the five cubic fields]
\label{corr:cubic-class-numbers}
Each of the five totally real cubic fields in
Table~\ref{corr:classnumber-table} has class number one.  In the table,
$\alpha$ denotes a root of the polynomial in its row, $d_H$ is the
field discriminant, and $B=\lfloor 2\sqrt{d_H}/9\rfloor$ is the
integer Minkowski bound.  The final column lists a generator of a
principal ideal of norm $p$ for each indicated rational prime $p$;
the field norm of the generator may be $-p$.
\end{proposition}

\begin{table}[htbp]
\centering
\small
\begin{tabular}{c|r|r|l}
 Polynomial for $\alpha$ & $d_H$ & $B$ & $p:\ \text{generator}$\\\hline
 $X^3-X^2-3X+1$ & $148$ & $2$ & $2:\alpha-1$\\
 $X^3-X^2-5X-1$ & $404$ & $4$ &
       $2:\alpha+1;\quad3:\alpha^2-2\alpha-2$\\
 $X^3-X^2-7X-3$ & $788$ & $6$ &
       $2:\alpha+1;\quad3:\alpha;\quad5:\alpha^2-2\alpha-4$\\
 $X^3-X^2-4X+3$ & $257$ & $3$ & $3:\alpha$\\
 $X^3-X^2-11X-7$ & $2708$ & $11$ &
       $2:\alpha+1;\quad3:\alpha+2;\quad7:\alpha$\\
 & & & $5:-\alpha^2+2\alpha+10;\quad11:2\alpha^2-4\alpha-17$
\end{tabular}
\caption{Exact generators used in the class-number-one proof.}
\label{corr:classnumber-table}
\end{table}

\begin{proof}
The rational-root test proves irreducibility of each polynomial $f$.
Their polynomial discriminants are, in the displayed order,
$148,404,788,257,2708$, all positive.  Thus all five fields are
totally real.  At discriminant $257$, squarefreeness immediately gives
$\mathcal O_{H_0}=\mathbb Z[\alpha]$.  In each other row only $2$
could divide the index.  In all four cases
$\overline f=(X+1)^3$ modulo $2$, while
\[
 \left.\frac{f(X)-(X+1)^3}{2}\right|_{X=1}
                  \in\{-5,-7,-9,-13\}.
\]
These values are odd.  Dedekind's index criterion therefore excludes
$2$ from the index.  Equivalently, each shifted polynomial $f(X+1)$
is Eisenstein at $2$, so its root generates the local ring of integers
\cite[Proposition~7.55 and Remark~7.56]{MilneANT}.  This proves
maximality and the asserted field discriminants.  The use of Dedekind factorization below is consequently
valid also at $2$.

Minkowski's theorem gives an integral ideal in every ideal class
with norm at most $2\sqrt{d_H}/9$ for a totally real cubic field
\cite{MilneANT}.  It suffices to prove principalness for every prime
ideal whose norm is at most $B$.  In the four even-discriminant rows,
$2$ is totally ramified; the displayed norm-$2$ element generates its
unique prime ideal.  At discriminant $257$, the polynomial is
irreducible modulo $2$, so the prime above $2$ has norm $8>B$.

For each relevant odd rational prime $p\leq B$, direct factorization
of $f$ modulo $p$ has type $1+2$.  There is therefore exactly one
prime ideal of norm $p$.  The corresponding element in the table
has norm $\pm p$, so this degree-one prime is principal.  Its
degree-two partner is principal as well, since their product is
$(p)$.  This last observation covers in particular the prime of
norm $9$ in the field of discriminant $2708$.  The table accounts
for all rational primes $p\leq B$, hence for all prime ideals within
the Minkowski bound.  Every ideal class is therefore trivial.

All finite calculations in this proof---discriminants, index tests,
modular splitting types, and generator norms---are reproduced with
integer arithmetic by \path{code/verify_cubic_class_numbers.py}.
\end{proof}

\subsection{Gamma relations and the scope of exact certificates}
\label{corr:gamma-certificates}

The intake README already identifies a major attribution correction
in \path{gamma-lattice-and-certificates.tex}.  Its \texttt{law:ko}
and corresponding abstract claim identify completeness of all
multiplicative relations among rational gamma values with a theorem
of Koblitz--Ogus.  The sufficient algebraicity criterion is a theorem;
the assertion that every algebraic gamma monomial follows from the
standard relations is Rohrlich's conjecture.  The distinction is
explicit, for example, in the gamma-value discussion of
Rivoal and Roques \cite{RivoalRoques2014}.

The exact reductions already obtained from reflection and Gauss
multiplication remain valid.  Their row reduction gives completeness
within the stated formal relation space.  Describing the remaining
numerical gamma values as irreducible among all possible algebraic
monomial relations requires the conjecture.  Accordingly,
\texttt{law:ko} should become a conjecture or a conditional statement,
and the same qualification must be propagated to the abstract,
introduction, and basis claims.  A negative integer-relation search
cannot supply the missing converse.

There is a separate issue in the certificate model.  The article
explains that its principal-branch functional-equation instances
are screened by a numerical branch-domain guard.  Exact cancellation
of all special-function atoms verifies the algebraic replay, but
proves the target identity only when the instantiated functional
identities and their branch conditions are valid.  A finite-precision
small-residual check is not a proof of those conditions.  A complete
certificate should therefore include an exact proof of each relevant
branch condition, or explicitly list that condition as an assumption.
For algebraic substitution points, exact real-algebraic sign tests
and a specified argument sector can often provide this missing part.
Single-valued Bloch--Wigner identities offer another route for claims
that can be expressed entirely in that language.

This qualification does not assert an error in any particular
certified value.  The checker and identity stores themselves were
not included in the inspected directory, so their actual soundness
cannot be audited from these reports alone.  It specifies the proof
obligation left by the model described in the article.

% ----- sections/corrections_additional.tex -----
\subsection{Clausen parity and other already identified source issues}
\label{audit:additional}

There is a direct convention mismatch in
\path{gaussian-eisenstein-double-polylogs.tex}, equation
\texttt{eq:eis-single}. With $\omega=e^{2\pi i/3}$, the identity
\[
 \Re\Li_n(\omega)=\tfrac12(3^{1-n}-1)\zeta(n),\qquad n>1,
\]
is correct, but the assertion
$\Im\Li_n(\omega)=\operatorname{Cl}_n(2\pi/3)$ for every $n$
conflicts with the parity convention used elsewhere in the corpus.
Under the convention in Section~\ref{intro:section}, this is correct
only for even $n$. At odd $n$, use the sine series
$\sum_{m\geq1}\sin(2\pi m/3)/m^n$ explicitly. For example,
\begin{equation}\label{audit:clausen-three}
 \Im\Li_3(\omega)=\frac{2\pi^3}{81},
 \qquad
 \operatorname{Cl}_3(2\pi/3)=\Re\Li_3(\omega)
                         =-\frac49\zeta(3).
\end{equation}
The first identity follows from the elementary Fourier series
\[
 \sum_{m\geq1}\frac{\sin(m\theta)}{m^3}
 =\frac{\pi^2\theta}{6}-\frac{\pi\theta^2}{4}
                         +\frac{\theta^3}{12},
 \qquad 0\leq\theta\leq2\pi,
\]
and the second from distribution. The two values even have opposite
signs. Numerical confirmation of a formula under an unstated
alternative convention cannot resolve this notation conflict.

The README already records a factor error in \path{report-1.tex}:
the displayed $\pi^2/6$ term in its questioned dilogarithm identity
must be $\pi^2/3$. This is an inherited correction, not a new finding
of the present audit. It should be applied together with regeneration
of the affected PDF. Likewise, the README's warnings about stale
Herglotz cross-references and the Koblitz--Ogus attribution should be
propagated to the source articles and their generated versions.

Finally, \path{stieltjes-parameter-derivative-tower.tex},
Section~\texttt{sec:atoms}, promotes unsuccessful integer-relation
searches for six second-$s$-derivative values to assertions of their
mutual independence. The reported searches support a list of
unreduced candidates for the tested coefficient bounds and inventory.
They do not prove the independence statement. The same qualification
applies to the CM-point non-reducibility claims in
\path{eisenstein-row-sums-cm-polygamma.tex} when their evidence is a
finite search. Our formal rank theorems do not supply the missing
arithmetic implication.

% ----- sections/verification.tex -----
\section{Computational verification and reproducible artifacts}
\label{verification:section}

The package separates exact finite certificates from numerical
diagnostics. Its proofs apply to unbounded derivative orders and
denominators. The programs verify finite presentations, constants,
signs, conventions, and selected asymptotic predictions; none uses
unsuccessful integer-relation searches as a proof of independence.
All reported runs completed successfully.

\subsection{Exact arithmetic checks}

\begin{itemize}
\item \path{code/verify_ranks.py} checks $486$ exact rational
presentations. The $54$ denominators are $1,\ldots,48$ and
$60,64,72,90,105,120$. The power weights use
$s=-4,-2,-1,0,1,2,3,5$; a ninth weight system includes zero and
non-power values. The checks compare prime and all-divisor row
spans, endpoint conventions, both reflection signs, and the
nonresonant unit-symbol basis.
\item \path{code/verify_jet_ranks.py} checks $160$ expanded finite-jet
presentations, including nilpotent weights. Exact fraction elimination
confirms the dimensions before and after all endpoint layers are
prescribed. The rational specializations test the algebraic theorem;
they do not approximate logarithms as rational numbers and then claim
to compute an arithmetic rank.
\item \path{code/verify_shuffle_audit.py} constructs the $13$ by $10$
shuffle-product matrix, verifies its rank $7$, and verifies the three
integer annihilators in Section~\ref{audit:poly:shuffle-rank}. It checks
that every individual coordinate is outside that row space.
\item \path{code/verify_cubic_class_numbers.py} uses only integer
arithmetic to check all finite premises of
Proposition~\ref{corr:cubic-class-numbers}: irreducibility,
discriminants, maximality tests, splitting types, and generator norms.
The deduction of class number one uses the stated Dedekind and
Minkowski theorems.
\end{itemize}

The JSON files in \path{data/} retain the exact parameters, rational
matrix information, and witnesses. In particular, the class-number
calculation includes degree-two prime ideals within the Minkowski
bound; checking only rational primes with degree-one generators
without accounting for their partners would leave a gap.

\subsection{Independent numerical representations}

The Stieltjes script uses $85$ decimal digits. It compares the
Laplace kernel with scaled Hurwitz-jet evaluation in $8$ cases,
computes $40$ moving zeros for $n=1,2,3,4$ and
$k=8,32,128,512$, and verifies the global first-index bracket at
$k=1,2,4,8,16,32,64,128$. The largest kernel discrepancy, divided
by the larger of $1$ and the comparison magnitude, was less than
$1.7\cdot10^{-86}$. This is a residual at the selected working
precision, not a certified bound on the error in every printed digit.

There is a practical numerical issue at large $k$. The raw value of
$\zeta(k+1,a)$ can be extremely small on the moving-zero scale
$a\asymp k$. An absolute-accuracy zeta routine may therefore return
an adequate unscaled answer that loses many relative digits when
multiplied by $a^{k+1}$. The supplied program avoids that loss by
scaling \emph{inside} Euler--Maclaurin summation. It evaluates
\[
 \frac{a^{k+1}}{k!}F_{n,k}(a)
 =n![z^n]\left\{
     \frac{(1+z)_k}{k!}\,a^{k+1}\zeta(k+1+z,a)
                  \right\}
\]
with truncated power series in $z$ and terms of the form
$\exp(-(k+1)\log(1+m/a))$. This avoids finite differencing in the
jet variable. Each computed root is checked again after adding
$31$ terms to the finite Euler--Maclaurin sum. The stopping rule is
numerical, and explicit interval tail enclosures remain a proposed
extension.

For $n=1$, Table~\ref{verification:zeros} illustrates the improvement
from the explicit $1/(k+1)$ coefficient in
Section~\ref{zero:asymptotic-section}. Let $A_k$ denote the
leading term plus constant offset, and put $K=k+1$.
Here the next coefficient is
$c=0.0288707254210824205588\ldots$.

\begin{table}[htbp]
\centering
\small
\begin{tabular}{r r r r}
\toprule
$k$ & computed $a_k$ & $|a_k-A_k|$ & $|a_k-A_k-c/K|$\\
\midrule
$32$ & $56.558158770512722$ & $9.2190\cdot10^{-4}$ & $4.7031\cdot10^{-5}$\\
$128$ & $227.540415772647917$ & $2.2678\cdot10^{-4}$ & $2.9726\cdot10^{-6}$\\
$512$ & $911.472053968737555$ & $5.6465\cdot10^{-5}$ & $1.8632\cdot10^{-7}$\\
\bottomrule
\end{tabular}
\caption{Numerical zero locations and errors of two proved asymptotic
approximations. The displayed digits are rounded diagnostics.}
\label{verification:zeros}
\end{table}

The bridge script checks $12$ cases at $100$ decimal digits, including
nonreal primitive characters. The largest absolute residual was
less than $2.9\cdot10^{-101}$. Deliberately omitting the factor
$1/2$ or the required conjugation produces detectable errors, so
these checks test precisely the delicate parts of the formula.

The Gaussian script checks $p=1,3,5,7,9$ in three ways: accelerated
series summation, quadrature, and the finite reduction.
At $90$ digits, their largest pairwise discrepancy was less than
$4\cdot10^{-91}$. It also checks five real or complex generating
function arguments and the resummation at $z=1$.
The correction scripts supply $11$ checks at $70$ digits for the
log-gamma and Herglotz identities, and $7$ checks at $80$ digits
for the modified-polylogarithm and root-of-unity witnesses.
The latter reproduce both the corrected and printed values of the
complex-embedding $P_4$ example.

\subsection{Build and integration}

The principal deliverables are
\path{polylogarithms_research.tex} and
\path{polylogarithms_research.pdf}. The main \LaTeX{} file is
assembled from the editable fragments in \path{sections/}; it
contains the full text and bibliography and depends only on the
included figure and ordinary \LaTeX{} packages.
\path{BUILD.sh} rebuilds it using \texttt{pdflatex}.
\path{code/run_checks.py} reruns the verification programs.
The exact Python package versions are recorded in
\path{requirements.txt} and \path{ENVIRONMENT.json}.

The correction register records proposed changes to existing files;
the package does not silently replace those source drafts.
\path{PROVENANCE.json} fixes the inspected repository state, and
\path{SHA256SUMS} records checksums of the delivered artifacts.
The figures and data are reproducible from the included programs.
This makes each proof, correction, and finite witness independently
reviewable during repository integration.

% ----- sections/research_questions.tex -----
\section{Further research questions}
\label{questions:section}

The following questions follow directly from proved results or identified
gaps.  They are not presented as solved by the computations in this package.

\subsection{Zero geometry and asymptotic analysis}

\paragraph{1. The first derivative order at which all zeros appear.}
Let $k_0(n)$ be the least $k\geq1$ for which $\gamma_n^{(k)}$ has
$n$ distinct positive zeros.  Our proof shows that $k_0(n)$ exists,
and interlacing shows that saturation persists at every later order.
Determine $k_0(n)$, or obtain an explicit upper bound.
Is $k_0(n)=1$ for every $n$?  This is a concrete question, not a
consequence of our large-$k$ argument.  A useful first project is a
rigorously enclosed zero table for moderate $n$ and all smaller $k$.

\paragraph{2. An effective saturation theorem.}
Turn the proof into a displayed bound $k\geq K(n)$ by bounding the root
separation of $P_n$, the derivatives of $h_n$ near its zeros, and the
gamma concentration remainder.  This would replace a fixed-$n$
existence result by an algorithm with a proved termination bound.
The finite-factor proof of simple roots is qualitative; quantitative
root separation is the main missing input.

\paragraph{3. Joint growth of Laurent index and derivative order.}
The remainder constants currently depend on $n$ without control.
Find a range $n=n(k)\to\infty$ in which the zero count is saturated and
the expansions remain uniform.  This requires a quantitative theory
of the reciprocal-gamma Appell roots, not only fixed-degree
hyperbolicity.  Compare the resulting root distribution with known
asymptotic theory of Jensen and Appell polynomials.

\paragraph{4. Optimal bounds for the first-index root.}
The width-$1/2$ bracket is global and elementary.  Determine the best
constant $C$ for $a_k<e^{H_{k-1}}+C$, or prove sharper two-sided
bounds involving one or two elementary functions of $k$.
The expansion suggests precise limiting offsets, but monotonicity of
those offsets still needs proof.

\paragraph{5. Certified computation of all moving zeros.}
Implement the kernel and the scaled Euler--Maclaurin scheme with ball
arithmetic and explicit tail enclosures.  Combine signs at endpoints
with the proved zero bound to certify completeness.  This would give
an auditable bridge from the analytic theorem to large tables,
including coefficients of high-order zero expansions.

\paragraph{6. Interlacing in the Laurent index.}
The polynomial roots interlace in $n$, and the transformed zeros
eventually inherit the corresponding leading-order ordering.
Determine whether an exact interlacing law holds for
$\gamma_n^{(k)}$ and $\gamma_{n+1}^{(k)}$ at a fixed $k$, and identify
any necessary threshold or counterexample.  Positivity of the kernel
alone does not settle this comparison between different polynomials.

\paragraph{7. Complex zeros and other asymptotic scales.}
Extend the concentration analysis to complex $a$ in sectors.
The real sign-variation bound has no direct complex analogue, so a
complex zero-completeness result requires a separate argument, for
example an argument-principle estimate or a uniform saddle analysis.
Potential transitions near poles of $1/(1-e^{-t})$ should be treated
explicitly rather than inferred from the real theorem.

\subsection{Relations, depth, and exact certificates}

\paragraph{8. Arithmetic injectivity of the formal quotient.}
Specify manageable coefficient fields and background inventories for
which the evaluation map from a weighted distribution quotient is
injective, or has a computable kernel.  The formal rank is now known;
the arithmetic obstruction has consequently been isolated precisely.
Unsuccessful PSLQ searches can suggest test cases but cannot settle
this injectivity problem.

\paragraph{9. Integral presentations and torsion.}
The character proof works over characteristic zero after scalar
extension.  Determine the Smith normal form of the corresponding
integral matrices for positive integer weights, and study denominators
and torsion for negative weights after clearing denominators.
Rational dimension and integral freeness are different questions;
the arbitrary-weight theorem here does not assert the latter.

\paragraph{10. The odd-total-weight Gaussian ladder.}
For $S_{2m}$, compute explicit motivic or coaction classes in a precisely
specified level-four algebra.  The even-weight upper reductions are
proved, but the proposed lower-depth obstruction in the complementary
parity still needs such a calculation.  Start with a small named family
and produce a checkable nonvanishing certificate before asserting a
general alternation theorem.

\paragraph{11. The three-dimensional shuffle quotient at weight six.}
Determine what happens to its three surviving formal directions after
specialization to Gaussian or Eisenstein arguments.  Additional
functional equations, distribution identities, or motivic relations
may lower the dimension.  The exact annihilator witnesses in the
package provide a reproducible starting point, not a proof that the
specialized periods are independent.

\paragraph{12. Complete correction and certificate migration.}
Recompute the modified-polylogarithm complex-embedding tables with the
correct summation endpoint; reconcile Clausen parity conventions;
replace numerical branch guards by exact domain certificates; and
extend the five class-number certificates proved here to the other
fields in subsequent Herglotz investigations. Proving the fitted
Herglotz regulator evaluations remains a separate task even in those
five class-number-one fields. These are well-specified projects
that would make the existing experimental corpus substantially easier
to trust and reuse.

\paragraph{13. Formalization of the elementary core.}
Formalize prime-row factorization, conductor transport, the polynomial
root-preserving lemma, and the Rolle induction for Laplace transforms
in Lean. These reusable elementary results provide manageable foundations
for a later formal treatment of the analytic tower.

% ----- sections/bibliography.tex -----
\clearpage
\phantomsection
\addcontentsline{toc}{section}{References}
\begingroup
\small
\setlength{\parskip}{0pt}
\begin{thebibliography}{99}
\setlength{\itemsep}{4pt}
\raggedright

\bibitem{RepoSnapshot}
V. Reshetnikov, \emph{ProveIt: Analysis/Polylogarithms},
repository snapshot
\texttt{3a6d80ed618d839915deb0d19ab687f45e81d995},
8 October 2026 UTC.
\url{https://github.com/VladimirReshetnikov/ProveIt/tree/3a6d80ed618d839915deb0d19ab687f45e81d995/Analysis/Polylogarithms/docs}.
The intake README is in the parent directory; source paths and Git
blob identifiers are listed in the accompanying \path{PROVENANCE.json}.

\bibitem{DraftTower}
ProveIt draft corpus, \emph{The Parameter-Derivative Tower of the
Generalized Stieltjes Constants: a Uniform Harmonic-Number Theorem}, source
\path{docs/articles/stieltjes-parameter-derivative-tower.tex},
in the snapshot \cite{RepoSnapshot}.

\bibitem{DraftLadder}
ProveIt draft corpus, \emph{Stieltjes Antiderivatives, Hurwitz Zeta Jets,
and Polygamma Functions of Negative Order}, source
\path{docs/articles/stieltjes-antiderivative-ladder.tex},
in the snapshot \cite{RepoSnapshot}.

\bibitem{DLMF25}
NIST, \emph{Digital Library of Mathematical Functions},
Chapter 25: Zeta and Related Functions, especially
\S25.11 (Hurwitz zeta) and \S25.15 (Dirichlet $L$-functions).
\url{https://dlmf.nist.gov/25.11};
\url{https://dlmf.nist.gov/25.15}.
Accessed 8 October 2026.

\bibitem{DLMF5}
NIST, \emph{Digital Library of Mathematical Functions},
Chapter 5: Gamma Function, especially \S5.4 (special values),
\S5.5 (functional relations), and \S5.8 (infinite products).
\url{https://dlmf.nist.gov/5.4};
\url{https://dlmf.nist.gov/5.5};
\url{https://dlmf.nist.gov/5.8}.
Accessed 8 October 2026.

\bibitem{Coffey2009}
M. W. Coffey, “Series representations for the Stieltjes constants,”
\emph{Rocky Mountain Journal of Mathematics} \textbf{44} (2014),
443--477; preprint arXiv:0905.1111 (2009).
\url{https://arxiv.org/abs/0905.1111}.
Proposition 5 treats the unique maximum of $\gamma_1(a)$.

\bibitem{Kubert1979}
D. S. Kubert, “The universal ordinary distribution,”
\emph{Bulletin de la Soci\'et\'e Math\'ematique de France}
\textbf{107} (1979), 179--202.
\url{https://www.numdam.org/item/BSMF_1979__107__179_0.pdf}.

\bibitem{gau-olaikhan-2022}
A. Olaikhan, answers on Mathematics Stack Exchange:
“How to write a limit in terms of finite summation,”
3 March 2022, edited 5 March 2022,
\url{https://math.stackexchange.com/questions/4139313/how-to-write-a-limit-in-terms-of-finite-summation/4395416};
and the all-weight generalization posted on 5 March 2022 to
“Evaluate $\sum_{n=1}^{\infty}(-1)^{n-1}H_n/(2n+1)^3$,”
\url{https://math.stackexchange.com/questions/3262663/evaluate-sum-n-1-infty-frac-1n-1h-n2n13}.

\bibitem{audit-glanois-2016}
C. Glanois, “Motivic unipotent fundamental groupoid of
$\mathbb G_m\setminus\mu_N$ for $N=2,3,4,6,8$ and Galois descents,”
\emph{Journal of Number Theory} \textbf{160} (2016), 334--384.
\href{https://doi.org/10.1016/j.jnt.2015.08.003}{doi:10.1016/j.jnt.2015.08.003};
\url{https://arxiv.org/abs/1411.4947}, especially \S3.1 of version 2.

\bibitem{audit-pari-polylog}
The PARI Group, \emph{PARI/GP documentation: Transcendental functions},
entry \texttt{polylog}, flag $3$ (Zagier's modified polylogarithm).
\url{https://pari.math.u-bordeaux.fr/dochtml/html-stable/Transcendental_functions.html}.
Accessed 8 October 2026.

\bibitem{EspinosaMoll2002b}
O. Espinosa and V. H. Moll, “On Some Integrals Involving the
Hurwitz Zeta Function: Part 2,” \emph{Ramanujan Journal}
\textbf{6} (2002), 449--468.
\url{https://www.math.tulane.edu/~vhm/papers_html/hurwitz2.pdf}.

\bibitem{RadchenkoZagier2023}
D. Radchenko and D. Zagier, “Arithmetic properties of the Herglotz
function,” \emph{Journal f\"ur die reine und angewandte Mathematik}
\textbf{797} (2023), 229--253.
\url{https://arxiv.org/abs/2012.15805}.

\bibitem{MilneCFT}
J. S. Milne, \emph{Class Field Theory}, version 4.03,
6 August 2020.
\url{https://www.jmilne.org/math/CourseNotes/CFT.pdf}.

\bibitem{MilneANT}
J. S. Milne, \emph{Algebraic Number Theory}, version 3.08,
19 July 2020, especially Theorem 4.3 and
Proposition 7.55 with Remark 7.56.
\url{https://www.jmilne.org/math/CourseNotes/ANT.pdf}.

\bibitem{RivoalRoques2014}
T. Rivoal and J. Roques, \emph{Hadamard products of algebraic functions},
author preprint, 29 April 2014.
\url{https://rivoal.perso.math.cnrs.fr/articles/algcoef.pdf}.
See the discussion of gamma relations and Rohrlich's conjecture.

\end{thebibliography}
\endgroup
\end{document}
